\section{VI. 법체계의 토대들(THE FOUNDATIONS OF A LEGAL
SYSTEM)}\label{vi.-uxbc95uxccb4uxacc4uxc758-uxd1a0uxb300uxb4e4the-foundations-of-a-legal-system}

\subsection{1. 승인 규칙과 법적 유효성(RULE OF RECOGNITION AND LEGAL
VALIDITY)}\label{uxc2b9uxc778-uxaddcuxce59uxacfc-uxbc95uxc801-uxc720uxd6a8uxc131rule-of-recognition-and-legal-validity}

제4장에서 비판된 이론에 따르면, 한 법체계의 토대들은 사회 집단의 다수가,
스스로는 누구에게도 습관적으로 복종하지 않는 주권자인 인격 또는 인격들의
위협으로 뒷받침된 명령들에 습관적으로 복종하는 상황으로 이루어진다. 이
이론에서 이러한 사회적 상황은 법의 존재에 대한 필요조건이자
충분조건이다. 우리는 이미 이 이론이 현대의 국내법 체계의 몇몇 두드러진
특징들을 설명할 수 없는 무능력(incapacity)을 상당히 상세하게 보여
주었다. 그럼에도 불구하고, 많은 사상가들의 정신을 사로잡아 온 이 이론이
시사하듯이, 그것은 비록 흐릿하고 오도적인 형태이기는 하나, 법의 몇몇
중요한 측면들에 관한 일정한 진리들을 포함하고 있기는 하다. 그러나 이러한
진리들은, 의무의 일차 규칙들을 식별하기(identify) 위해 이차적 승인
규칙(rule of recognition)이 수용되고 사용되는 보다 복잡한 사회적 상황의
관점에서만 비로소 명확하게 제시될 수 있고, 그 중요성 또한 올바르게
평가될 수 있다. 바로 이 상황이야말로, 무엇인가가 한 법체계의
토대들이라고 불릴 만하다면, 그렇게 불릴 자격이 있는 것이다. 이 장에서는
주권 이론 및 다른 곳들에서 부분적으로만 또는 오도된 방식으로 표현되어 온
이 상황의 여러 요소들을 논의할 것이다.

이러한 승인 규칙이 수용되는 곳에서는 어디에서나, 사적 개인들과
공무담당자들 모두에게 의무의 일차 규칙들을 식별하기(identify) 위한 권위
있는 기준들이 제공된다. 앞서 보았듯이, 이렇게 제공되는 기준들은 다음의
다양한 형태들 가운데 하나 또는 그 이상의 형태를 취할 수 있다: 여기에는
권위 있는 문헌에 대한, 입법적 제정에 대한, 관습적 실천에 대한, 특정된
사람들의 일반적 선언에 대한, 또는 개별 사건들에서의 과거 사법적 결정들에
대한 참조(reference)가 포함된다. 제4장에서 묘사된 Rex I의 세계와 같이
매우 단순한 체계에서는, 그가 제정한 것만이 법이며 관습 규칙이나 헌법적
\textbf{(p.~101)} 문헌에 의해 그의 입법 권한에 법적 제한들이 전혀
부과되지 않기 때문에, 법을 식별하는 유일한 기준은 Rex I에 의한
제정이라는 사실에 대한 단순한 참조가 될 것이다. 이러한 단순한 형태의
승인 규칙의 존재는, 공무담당자들이나 사적 개인들이 이 기준에 따라
규칙들을 식별하는 일반적 실천 속에서 분명히 드러난다. 법의 다양한
‘원천들(sources)’이 존재하는 현대의 법체계에서는 승인 규칙도 이에
상응하여 더 복잡해지며, 법을 식별하기 위한 기준들은 복수적이고,
통상적으로 서면 헌법, 입법부에 의한 제정, 그리고 사법적 선례를 포함한다.
대부분의 경우 이러한 기준들 사이에 발생할 수 있는 충돌을 대비하여,
그것들을 상대적 종속성과 우위성(primacy)의 질서 속에 배열하는 규정이
마련되어 있다. 바로 이러한 방식으로, 우리 체계에서는 ‘보통법(common
law)’이 ‘제정법(statute)’에 종속된다.

한 기준이 다른 기준에 대해 갖는 이러한 상대적
\emphksans{종속성(subordination)}은 \emphksans{도출(derivation)}과 구별하는 것이
중요한데, 왜냐하면 이 두 관념을 혼동함으로써, 모든 법은 본질적으로 또는
‘실제로(really)’(설령 단지 ‘묵시적으로(tacitly)’일 뿐이라 하더라도)
입법의 산물이라는 견해에 대해 그럴듯하지만 허위적인 지지가 제공되어 왔기
때문이다. 우리 체계에서 관습과 판례는 입법에 종속되는데, 이는 관습법 및
보통법 규칙들이 제정법(statute)에 의해 그 법적 지위를 박탈당할 수 있기
때문이다. 그러나 이 규칙들은, 그 지위가 불안정하다는 점은 인정하더라도,
법으로서의 지위를 ‘묵시적인’ 입법 권한의 행사에 빚지고 있는 것이 아니라,
그것들에게 이러한 독립적이되 종속된(subordinate) 지위를 부여하는 승인
규칙의 수용에 빚지고 있다. 또, 단순한 경우에서와 마찬가지로, 이러한
상이한 기준들이 위계적으로 배열된 복합적인 승인 규칙의 존재는, 그러한
기준들에 의해 규칙들을 식별하는 일반적 실천 속에서 드러난다.

법체계의 일상적 운영 속에서 승인 규칙(rule of recognition)은 규칙으로서
명시적으로 정식화되는(expressly formulated) 경우가 극히 드물다. 다만
영국의 법원들이 때때로 법의 한 기준이 다른 기준에 대해 차지하는 상대적
지위를 일반적인 표현으로 공표하는(announce) 경우는 있는데, 예컨대 의회
제정법들(Acts of Parliament)이 법의 다른 원천들(sources)이나 법의 후보
원천들(suggested sources of law)에 대해 최고성(supremacy)을 가진다고
단언하는(assert) 때가 그러하다. 대체로 승인 규칙은 진술되지 않지만,
법원이나 다른 공무담당자들, 또는 사적 개인들이나 그들의 조언자들이 개별
규칙들을 식별하는(identify) 방식 속에서 그 존재가
\emphksans{드러난다(shown)}. 물론 승인 규칙이 제공하는 기준들을 법원이
사용하는 방식과, 다른 이들이 사용하는 방식 사이에는 차이가 있다.
왜냐하면 법원이 특정한 규칙이 법으로서 올바르게 식별되었다는 전제하에
어떤 특정한 \textbf{(p.~102)} 결론에 도달할 경우, 그들이 하는 말은 다른
규칙들에 의해 부여된 특별한 권위 있는 지위를 가지기 때문이다. 이러한
점에서, 그리고 다른 많은 점에서, 법체계의 승인 규칙은 게임의 득점
규칙(scoring rule)과 유사하다. 게임의 진행 과정에서 득점을 구성하는
행위들(득점(runs), 골 등)을 규정하는 일반 규칙은 거의 정식화되지 않으며,
대신 경기 임원들과 선수들이 승리에 산입되는 특정 국면들을 식별하는 데
그것을 \emphksans{사용(use)}한다. 여기에서도 경기 임원들(심판이나
득점기록원)의 선언은 다른 규칙들에 의해 부여된 특별한 권위 있는 지위를
가진다. 더 나아가, 두 경우 모두에서 이러한 규칙의 권위 있는 적용들과, 그
규칙의 문언에 비추어 볼 때 명백히 요구되는 바에 대한 일반적 이해 사이에
충돌이 발생할 가능성이 존재한다. 이는 우리가 뒤에서 보게 되겠지만,
이러한 유형의 규칙 체계가 존재한다는 것이 무엇을 의미하는지에 대한
어떠한 설명에서도 반영되어야 할 복잡성이다.

법원과 그 밖의 사람들에 의해, 체계의 개별 규칙들을 식별하는(identify)
데에 명시되지 않은 승인 규칙(rule of recognition)이 사용되는 것은 내부
관점(internal point of view)의 특징이다. 이러한 방식으로 승인 규칙을
사용하는 사람들은, 그것들을 향도규칙들(guiding rules)로 자신들이
수용하고 있음을 드러내며, 이러한 태도와 함께 외부 관점(external point of
view)의 자연스러운 표현들과는 다른, 특징적인 어휘가 수반된다. 아마도
이러한 표현들 가운데 가장 단순한 것은 ‘\ldots 은 법이다(It is the law
that \ldots)’라는 표현일 것인데, 이는 판사들뿐만 아니라 법체계 아래에서
살아가는 보통 사람들의 입에서도, 체계의 어떤 주어진 규칙을 식별할 때
발견될 수 있다. 이는 ‘아웃(Out)’이나 ‘골(Goal)’이라는 표현과 마찬가지로,
다른 사람들과 함께 그 목적에 적절한 것으로 인정하는(acknowledges)
규칙들에 대한 참조(reference)로 한 사람이 상황을 평가하는(assessing a
situation) 언어이다. 이러한 규칙에 대한 공유된 수용(shared acceptance)의
태도는, 사회 집단이 그러한 규칙들을 수용하고 있다는 사실을
\emphksans{외부로부터(ab extra)} 기록할 뿐, 자신은 그 규칙들을 수용하지 않는
관찰자의 태도와 대비된다. 이 외부 관점의 자연스러운 표현은 ‘\ldots 은
법이다(It is the law that \ldots)’가 아니라, ‘영국에서는 의회 내 여왕이
제정하는 것은 무엇이든 법으로 인정한다(In England they recognize as law
\ldots{} whatever the Queen in Parliament enacts\ldots)’와 같은 것이다.
이 두 표현 형식 가운데 첫 번째를 우리는 \emphksans{내부 진술(internal
statement)}이라고 부를 것인데, 이는 내부 관점을 드러내며, 승인 규칙을
수용하고 있으면서도 그것이 수용되고 있다는 사실을 진술하지 않은 채, 그
규칙을 적용하여 체계의 어떤 특정 규칙을 유효한(valid) 것으로 인정하는
사람이 자연스럽게 사용하는 표현이기 때문이다. \textbf{(p.~103)} 두 번째
표현 형식은 \emphksans{외부 진술(external statement)}이라고 부를 것인데, 이는
체계의 승인 규칙을 스스로 수용하지 않으면서, 다른 사람들이 그것을
수용하고 있다는 사실을 진술하는 외부 관찰자의 자연스러운 언어이기
때문이다.

수용된 승인 규칙을 사용하여 내부 진술을 하는 이러한 방식이 이해되고, 그
규칙이 수용되고 있다는 사실에 대한 외부 사실 진술(external statement of
fact)과 신중하게 구별된다면, 법적 ‘유효성(validity)’이라는
관념(notion)을 둘러싼 많은 모호함은 사라진다. 왜냐하면
‘유효한(valid)’이라는 말은 언제나 그런 것은 아니지만, 가장 흔히는 바로
그러한 내부 진술에서 사용되며, 법체계의 어떤 특정 규칙에 대하여,
명시되지는 않았으나 수용된 승인 규칙을 적용하기 때문이다. 어떤 주어진
규칙이 유효하다고 말하는 것은, 그것이 승인 규칙이 제공하는 모든 검사
기준들(tests)을 통과하였음을, 따라서 체계의 규칙임을 인정하는 것이다.
기실 우리는, 어떤 특정 규칙이 유효하다는 진술이란 그것이 승인 규칙이
제공하는 모든 기준들(criteria)을 충족한다는 것을 의미한다고 간단히 말할
수 있다. 이는 그러한 진술들의 내부 성격을 가릴 수 있다는 점에서만
부정확한데, 왜냐하면 이러한 유효성 진술들은 크리켓 선수의
‘아웃(Out)’이라는 표현과 마찬가지로, 화자와 다른 사람들이 수용하고 있는
승인 규칙을 어떤 특정 사례에 적용하는 것이지, 그 규칙이 충족되었다는
사실을 명시적으로 진술하는 것이 아니기 때문이다.

법적 유효성(legal validity)이라는 관념과 관련된 몇몇 난점들은, 법의
유효성과 법의 ‘실효성(efficacy)’ 사이의 관계에 관한 것이라고들 말해진다.
만약 ‘실효성’이란, 특정한 행위를 요구하는 어떤 법규칙이 대체로(more
often than not) 복종된다는 사실을 의미하는 것이라면, 개별 규칙의
유효성과 \emphksans{그 규칙의(its)} 실효성 사이에는 필연적인 연관(necessary
connection)이 없다는 점은 분명하다. 다만 체계의 승인 규칙(rule of
recognition)이 그 기준들(criteria) 가운데에---일부 체계가
그러하듯---어떤 규칙이 오랫동안 실효성을 상실한 경우에는 더 이상 그
체계의 규칙으로 간주되지 않는다는 규정(때때로 실효상실 규칙(rule of
obsolescence)이라 불리는 것)을 포함하는 경우는 예외이다.

어떤 개별 규칙의 실효성 결여(inefficacy)는---그것이 그 규칙의 유효성에
불리하게 작용할 수도 있고 그렇지 않을 수도 있지만---체계 전체의 규칙들에
대한 일반적인 무시(general disregard)와는 구별되어야 한다. 이러한 무시는
그 성격상 완전하고 또한 장기간 지속될 수 있는데, 그 경우 우리는 새로운
체계에 대해서는 그것이 특정 집단의 법체계로서 결코 확립된 적이 없다고
말하게 되거나, 한때 확립되었던 체계에 대해서는 그것이 더 이상 그 집단의
법체계가 아니게 되었다고 말하게 될 것이다. 어느 경우든, 체계의 규칙들에
근거하여 \textbf{(p.~104)} 내부 진술(internal statement)을 하는 데
필요한 정상적인 맥락 또는 배경은 결여되어 있다. 이러한 경우에는, 특정
개인들의 권리들과 책무들을 그 체계의 일차 규칙들을 참조하여 평가하는
것도, 그 체계의 승인 규칙들을 참조하여 그 규칙들 가운데 어느 것의
유효성을 평가하는 것도, 일반적으로 \emphksans{무익한(pointless)} 일이 된다.
실제로 한 번도 실효적이었던 적이 없거나 이미 폐기된 규칙 체계를
적용하겠다고 고집하는 것은, 아래에서 언급할 특수한 사정을 제외하면, 한
번도 수용된 적이 없거나 이미 폐기된 득점 규칙(scoring rule)을 참조하여
경기의 진행 상황을 평가하려는 것만큼이나 무익하다.

체계의 특정 규칙의 유효성에 관하여 내부 진술을 하는 사람은, 그 체계가
일반적으로 실효성을 갖고 있다는 외부 사실 진술의 참임(truth)을
\emphksans{전제(presuppose)}하고 있다고 말할 수 있다. 왜냐하면 내부 진술의
정상적인 사용은 그러한 일반적 실효성의 맥락 속에서 이루어지기 때문이다.
그러나 유효성에 관한 진술이 체계가 일반적으로 실효적이라는 것을
‘의미한다(mean)’고 말하는 것은 잘못일 것이다. 어떤 체계가 한 번도 확립된
적이 없거나 이미 폐기된 경우, 그 체계의 규칙의 유효성에 대해 말하는 것은
대체로 무익하거나 공허한(normally pointless or idle) 말이기는 하지만,
그렇다고 해서 그것이 전혀 의미 없거나(meaningless) 언제나 무익한(always
pointless) 것은 아니다. 로마법(Roman Law)을 가르치는 하나의 생생한
방식은, 그 체계가 여전히 실효적\emphksans{인 것처럼(as if)} 말하며, 개별
규칙들의 유효성을 논하고 그 규칙들에 따라 문제를 해결하는 것이다; 또한
혁명에 의해 파괴된 옛 사회질서의 회복에 대한 희망을 품고 새로운 질서를
거부하는 하나의 방식은, 구체제의 법적 유효성 기준들(criteria of legal
validity)에 집착하는 것이다. 이는 차르 체제 러시아(Tsarist Russia)에서
유효했던 어떤 상속 규칙(rule of descent)에 근거하여 여전히 재산권을
주장하는 백계 러시아인(White Russian)에 의해 암묵적으로 행하여진다.

체계의 특정 규칙이 유효하다는 내부 진술과, 그 체계가 일반적으로
실효적이라는 외부 사실 진술 사이의 정상적인 맥락적 연관을 파악하면,
규칙의 유효성을 단언하는(assert) 것은 그것이 법원에 의해 집행되거나 그
밖의 어떤 공무담당 작용이 취해질 것임을 예측하는 것이라는 통상적 이론을
그에 걸맞은 관점에서 이해하는 데 도움이 된다. 이러한 이론은 우리가 지난
장에서 검토하고 배척한 의무의 예측적 분석과 여러 면에서 유사하다. 두
경우 모두에서 이 예측적 이론을 제시하게 되는 동기는, 오직 이러한
방식으로만 형이상학적 해석을 회피할 수 있다는 다음과 같은 확신이다: 즉,
규칙이 유효하다는 진술은 \textbf{(p.~105)} 경험적 수단으로는 탐지될 수
없는 어떤 신비한 성질을 귀속시키는 것이거나, 아니면 공무담당자들의 장래
행태에 대한 예측이어야 한다는 것이다. 또한 두 경우 모두에서 이 이론의
그럴듯함은 동일한 다음의 중요한 사실에 기인한다: 즉 체계가 일반적으로
실효적이며 앞으로도 그러할 개연성이 높다는, 관찰자가 기록할 수 있는,
외부 사실 진술의 참임(truth)이, 규칙을 수용하고 의무나 유효성에 관한
내부 진술을 하는 사람에 의해 정상적으로(normally) 전제된다는 점이다. 이
둘은 분명히 매우 밀접하게 연관되어 있다. 마지막으로, 두 경우 모두에서 이
이론의 오류는 동일하다. 그것은 내부 진술의 특수한 성격을 간과하고, 이를
공무담당 작용(official action)에 관한 외부 진술로 취급하는 데에 있다.

이러한 오류는 특정 규칙이 유효하다는 판사 자신의 진술이 사법적 결정에서
어떻게 기능하는지를 고려하면 즉각적으로 드러난다. 왜냐하면 이 경우에도
판사는 그러한 진술을 하면서 체계의 일반적 실효성을 전제할 뿐 이를
진술하지는 않지만, 분명히 자기 자신의 또는 타인의 공무담당 작용을
예측하는 데에는 관심이 없기 때문이다. 규칙이 유효하다는 그의 진술은, 그
규칙이 그의 법정에서 무엇이 법으로 간주될 것인지를 식별하기 위한 검사
기준들(tests)을 충족한다는 점을 인정하는 내부 진술이며, 결정에 대한
예언(prophecy)이 아니라 결정의 \emphksans{이유(reason)}의 일부를 구성한다.
기실 이러한 진술이 사적 개인에 의해 이루어지는 경우에는, 규칙이
유효하다는 진술을 예측으로 간주하는 것이 더 그럴듯해 보일 수 있다.
왜냐하면 비공무담당자의 유효성 또는 유효하지 않음에 관한 진술들과 사안을
판단하는 법원의 진술 사이에 충돌이 발생하는 경우, 전자는 철회되어야
한다고 말하는 데에 종종 상당한 일리가 있기 때문이다. 그러나 이 경우에도,
제7장에서 공무담당자들의 선언과 규칙의 명백한 요구 사이의 이러한 충돌의
의미를 검토하게 되면 보게 되듯이, 그것이 법원이 무엇이라고 말할지를
\emphksans{예측하는(predicted)} 데에서 오류를 범했기 때문에 이제
\emphksans{그른(wrong)} 진술로 드러나 철회되는 것이라고 단정하는 것은
독단적일 수 있다. 왜냐하면 진술을 철회하는 이유는 그것이 그르다는
사실만이 아니며, 또한 그른 방식 역시 이러한 설명이 허용하는 것보다 더
다양하기 때문이다.

다른 규칙들의 유효성을 평가하기 위한 기준들을 제공하는 승인 규칙(rule of
recognition)은 중요한 의미에서 \emphksans{궁극적(ultimate)} 규칙이다. 이 점은
우리가 곧 명확히 해설하려 한다. 그리고 통상 그렇듯이 여러 기준들이
상대적 종속성과 우위성(primacy)의 서열로 정렬되어 있는 경우, 그중 하나는
\emphksans{최고(supreme)} 기준이 된다. \textbf{(p.~106)} 승인 규칙의
궁극성(ultimacy)과 그 기준 중 하나의 최고성(supremacy)에 관한 이러한
관념들은 주의 깊게 다룰 만한 가치가 있다. 우리는 이 관념들을, 앞서
우리가 기각했던 이론---즉, 모든 법체계에는, 설령 그것이 법적 형식 배후에
숨어 있다 하더라도, 반드시 법적으로 무제한인 주권적 입법 권력(sovereign
legislative power)이 존재해야 한다는 이론---으로부터 분리해내는 것이
중요하다.

이 두 관념, 즉 최고 기준(supreme criterion)과 궁극적 규칙(ultimate rule)
가운데 전자는 정의하기가 가장 쉽다. 법적 유효성의 기준 또는 법의 원천이
최고(supreme)라고 말할 수 있는 경우는, 그 기준을 참조(reference)하여
식별된 규칙들이 다른 기준들을 참조하여 식별된 규칙들과 충돌하더라도
여전히 체계의 규칙으로 인정되는 반면, 후자의 기준들을 참조하여 식별된
규칙들은 최고 기준을 참조하여 식별된 규칙들과 충돌하는 경우 그러한
인정이 이루어지지 않을 때이다. 우리가 이미 사용해 온 ‘우위의(superior)’
기준과 ‘종속된(subordinate)’ 기준이라는 관념들도 이와 유사한 비교적
설명에 의해 해명될 수 있다. 우위의 기준과 최고 기준이라는 관념들은 단지
척도상에서의 \emphksans{상대적인(relative)} 위치를 가리킬 뿐이며, 법적으로
\emphksans{무제한인(unlimited)} 입법 권한에 대한 어떠한 관념도 함의하지
않는다는 점은 명백하다. 그럼에도 불구하고 ‘최고(supreme)’라는 것과
‘무제한적(unlimited)’이라는 것은---적어도 법이론에서는---쉽게 혼동된다.
그 이유 중 하나는, 보다 단순한 형태의 법체계에서는 궁극적 승인
규칙(ultimate rule of recognition), 최고 기준, 그리고 법적으로
무제한적인 입법부라는 관념들이 서로 수렴하는 것처럼 보이기 때문이다.
헌법적 제한을 받지 않는 입법부가 존재하고, 그 입법부가 자신의
제정(enactment)을 통해 다른 원천들로부터 나온 다른 모든 법규칙들의 법적
지위를 박탈할 권능을 가진(competent) 경우, 그러한 체계의 승인 규칙의
일부로서 그 입법부의 제정이 유효성의 최고 기준이 된다. 헌법이론에
따르면, 이것이 영국(United Kingdom)의 경우이다. 그러나 미합중국과 같이
그러한 법적으로 무제한적인 입법부가 존재하지 않는 체계들조차도, 유효성의
기준들로 이루어진 하나의 집합을 제공하는 궁극적 승인 규칙을 충분히
포함할 수 있으며, 그 기준들 가운데 하나는 최고일 수 있다. 이는 일반
입법부의 입법 권능(legislative competence)이 개정 권한을 포함하지
않거나, 또는 그 개정 권한의 범위 밖에 놓인 조항들을 포함하는 헌법에 의해
제한되는 경우에 해당한다. 이 경우, ‘입법부’라는 개념을 가장 넓게
해석하더라도 법적으로 무제한적인 입법부는 존재하지 않지만, 그 체계는
물론 궁극적 승인 규칙을 포함하고 있으며, 헌법 조항들 속에 유효성의 최고
기준을 포함하고 있다.

\textbf{(p.~107)} 승인 규칙(rule of recognition)이 한 체계의
\emphksans{궁극적(ultimate)} 규칙이라는 의미는, 매우 익숙한 하나의 법적 추론
연쇄를 따라가 보면 가장 잘 이해된다. 어떤 후보 규칙(suggested rule)이
법적으로 유효한지 여부가 문제로 제기되는 경우, 그 질문에 답하기 위해서는
다른 어떤 규칙이 제공하는 유효성의 기준을 사용해야 한다. 옥스퍼드셔
카운티 의회(Oxfordshire County Council)의 이 이른바 조례(purported
by-law)는 유효한가? 그렇다: 왜냐하면 그것은 보건부 장관(Minister of
Health)이 제정한 법정 명령(statutory order)에 의해 부여된 권한을
행사하여, 그 명령에 명시된 절차에 따라 제정되었기 때문이다. 이 첫
단계에서 법정 명령은 조례의 유효성이 평가되는 기준을 제공한다.
실천적으로는 더 나아갈 필요가 없을 수도 있지만, 그렇게 할 수 있는
상존하는 가능성은 있다. 우리는 그 법정 명령의 유효성을 문제 삼고, 그러한
명령을 제정할 권한을 장관에게 부여한 제정법(statute)을 참조하여 그
유효성을 평가할 수 있다. 마지막으로, 그 제정법의 유효성마저 문제 삼아
의회 내 여왕이 제정하는 것은 법이라는 규칙을 참조하여 평가하게 되면,
유효성에 관한 탐구는 여기서 멈추게 된다: 왜냐하면 이는 우리가 다른
규칙들의 유효성을 평가하기 위한 기준을 제공하는 규칙에 도달했기
때문이며, 동시에 그 규칙은 중간 단계의 법정 명령이나 제정법과는 달리, 그
자신의 법적 유효성을 평가하기 위한 기준을 제공하는 어떠한 규칙도 더 이상
존재하지 않는다는 점에서 그들과 구별된다.

기실 이 궁극적 규칙에 대하여 제기할 수 있는 질문들은 많다. 우리는
영국에서 법원들, 입법부들, 공무담당자들, 또는 사적 시민들이 실제로 이
규칙을 궁극적 승인 규칙으로 사용하고 있는 실천이 존재하는지 물을 수
있다. 아니면 우리의 법적 추론 과정은 이제 폐기된 체계의 유효성 기준을
가지고 벌이는 공허한 놀이에 불과한 것인가? 우리는 또한 이러한 규칙을
근간으로 하는 법체계가 만족스러운 형태의 법체계인지 물을 수 있다. 그것은
악보다 선을 더 많이 산출하는가? 그것을 지지할 타산적 이유들(prudential
reasons)은 존재하는가? 그렇게 할 도덕적 의무(moral obligation)는 있는가?
이러한 질문들은 분명히 매우 중요하다. 그러나 이와 동일하게 분명한 것은,
우리가 승인 규칙에 관하여 이러한 질문들을 제기할 때, 우리는 더 이상 그
규칙을 도구로 삼아 다른 규칙들에 관해 답했던 것과 같은 종류의 질문에
답하려 하고 있는 것이 아니라는 점이다. 우리가 어떤 특정한
제정(enactment)이 의회 내 여왕이 제정하는 것은 법이다라는 규칙을
충족하기 때문에 유효하다고 말하는 단계에서, 영국에서 바로 이 마지막
규칙이 법원들, 공무담당자들, 그리고 사적 개인들에 의해 궁극적 승인
규칙으로 사용되고 있다고 말하는 단계로 이동할 때, \textbf{(p.~108)}
우리는 체계의 한 규칙의 유효성을 단언하는 내부 법 진술(internal
statement of law)에서, 그 규칙을 수용하지 않더라도 체계의 관찰자가 할 수
있는 사실에 관한 외부 진술(external statement of fact)로 이동한 것이다.
마찬가지로, 어떤 특정한 제정이 유효하다는 진술에서, 그 법체계의 승인
규칙이 훌륭한 것이며 그것에 기초한 체계가 지지할 가치가 있는 체계라는
진술로 이동할 때, 우리는 법적 유효성에 관한 진술에서 가치에 관한 진술로
이동한 것이다.

일부 저자들은 승인 규칙(rule of recognition)의 법적 궁극성(legal
ultimacy)을 강조하면서, 체계의 다른 규칙들의 법적 유효성은 이를 참조하여
입증될(demonstrated) 수 있는 반면, 승인 규칙 자체의 유효성은 입증될 수
없고 ‘가정된다(assumed)’거나 ‘상정된다(postulated)’거나 하나의
‘가설(hypothesis)’이라고 표현해 왔다. 그러나 이러한 표현은 심각하게
오해를 불러일으킬 수 있다. 법체계의 일상적 삶 속에서 판사, 법률가, 또는
일반 시민에 의해 특정 규칙들에 관하여 행해지는 법적 유효성에 대한
진술들은 기실 일정한 전제들을 수반한다. 이러한 진술들은 체계의 승인
규칙을 수용하는 자들의 관점을 표현하는 내부 법 진술들(internal
statements of law)이며, 그러한 것으로서 체계에 관한 외부 사실
진술들(external statements of fact)에서 진술될 수 있는 많은 내용들을
명시하지 않은 채로 남겨 둔다. 이렇게 명시되지 않은 것들은 법적 유효성
진술들의 정상적인 배경 또는 맥락을 이루며, 따라서 그러한 진술들에 의해
‘전제된다(presupposed)’고 말해진다. 그러나 이러한 전제된 사항들이 정확히
무엇인지를 분명히 파악하고, 그 성격을 흐리게 하지 않는 것이 중요하다.
그것들은 두 가지로 구성된다. 첫째, 어떤 주어진 법규칙, 예컨대 특정한
제정법(statute)의 유효성을 진지하게 단언하는 사람은, 법을 식별하는 데
적절한 것으로 자신이 수용하는 승인 규칙을 스스로 사용하고 있다. 둘째,
그가 특정한 제정법의 유효성을 평가하는 기준으로 삼는 이 승인 규칙은 그
자신에게만 수용되는 것이 아니라, 체계의 일반적 작동 속에서 실제로
수용되고 사용되고 있는 승인 규칙이라는 점이다. 만약 이
전제(presupposition)의 참임이 의심된다면, 그것은 실제 실천을
참조함으로써 확립될 수 있다: 즉, 법원들이 무엇을 법으로 간주할지를
식별하는 방식, 그리고 이러한 식별들에 대한 일반적인 수용 또는
묵인(acquiescence)을 참조함으로써이다.

이 두 전제(presuppositions) 중 어느 것도, 입증할 수 없는
‘유효성(validity)’에 대한 ‘가정(assumption)’으로 제대로 설명될 수는
없다. 우리는 ‘유효성’이라는 단어를 오직 \textbf{(p.~109)} 규칙 체계
\emphksans{내부(within)}에서 어떤 규칙이 그 체계의 구성원으로서의 지위를 갖기
위해 승인 규칙(rule of recognition)에 의해 제공된 특정 기준들을 충족해야
할 때, 그러한 질문에 답하기 위해서만 필요로 하며, 대체로 그러한 경우에만
사용한다. 승인 규칙 자체의 유효성에 대해서는 이러한 질문이 제기될 수
없다. 그것은 유효하거나 유효하지 않은 것이 아니라, 단지 그러한 방식으로
사용되기에 적절한 것으로 수용될 뿐이다. 이러한 단순한 사실을, 그 규칙의
유효성은 ‘가정되었지만 입증될 수 없다(assumed but cannot be
demonstrated)’고 모호하게 표현하는 것은, 미터 단위 측정의
정확성(correctness)에 대한 궁극적 검사 기준(ultimate test)인 파리의 표준
미터 자(metre bar) 자체가 정확하다는 것을 ‘우리가 가정하지만 결코 입증할
수 없다’고 말하는 것과 같다.

보다 심각한 반론은, 궁극적 승인 규칙(ultimate rule of recognition)이
유효하다는 ‘가정(assumption)’에 관한 말이, 법률가들의 유효성 진술 배후에
놓인 두 번째 전제의 본질적으로 사실적(factual) 성격을 은폐한다는 점이다.
의심할 여지 없이, 승인 규칙의 실제 존재를 구성하는 판사들, 공무담당자들,
기타 행위자들의 실천(practice)은 복합적인 사안이다. 우리가 뒤에서 보게
되겠지만, 이러한 종류의 규칙의 정확한 내용과 범위, 나아가 그 존재 여부에
관한 질문들이 명확하거나 확정적인(determinate) 답변을 허용하지 않는
상황들도 분명히 존재한다. 그럼에도 불구하고, 이러한 규칙의 ‘유효성을
가정하는 것(assuming the validity)’과 그 ‘존재를 전제하는
것(presupposing the existence)’을 구별하는 것은 중요하다. 최소한 이러한
구별의 실패가, 그러한 규칙이 \emphksans{존재한다(exists)}고 단언하는 것이
무엇을 의미하는지를 흐리게 만들기 때문이다.

지난 장에서 개괄한 의무의 일차 규칙(primary rules of obligation)으로
이루어진 단순한 체계에서는, 어떤 특정 규칙이 존재한다고 단언하는 것은,
그 규칙을 수용하지 않는 관찰자가 제시할 수 있는 외부 사실 진술(external
statement of fact)에 불과하다. 이러한 진술은, 실제로 어떤 행위 양식이
일반적으로 표준으로 수용되고 있는지, 그리고 우리가 보았듯이 단순한
수렴적 습관들(convergent habits)과 사회적 규칙(social rule)을 구별하는
그러한 특징들이 수반되는지를 확인함으로써 검증될 수 있다. 마찬가지로,
현재 우리가 “잉글랜드에서는---비록 법적 규칙은 아니지만---교회에 들어갈
때 모자를 벗어야 한다는 규칙이 존재한다”는 단언을 해석하고 검증해야
하는 방식도 바로 이러한 방식이다. 이러한 유형의 규칙들이 사회 집단의
실제 실천 속에서 존재하는 것으로 확인된다면, 그 유효성에 관하여 별도로
논의할 문제는 없다. 물론 그 가치나 바람직성은 문제 삼을 수 있다. 그러나
일단 그 존재가 사실로 확립되었다면, 그것들이 유효하다고 말하거나
유효하지 않다고 부정하거나, \textbf{(p.~110)} 혹은 “우리는 그것들의
유효성을 가정했을 뿐 이를 입증할 수는 없다”고 말하는 것은 오히려 사안을
혼란스럽게 만들 뿐이다. 반면에, 성숙한 법체계와 같이 승인 규칙을
포함하는 규칙 체계를 갖는 경우에는, 어떤 규칙이 체계의 구성원으로서의
지위를 갖는지는 이제 승인 규칙이 제공하는 일정한 기준들을 충족하는지
여부에 달려 있다. 이는 ‘존재한다(exist)’라는 단어에 새로운 용법을
가져온다. 이제 어떤 규칙이 존재한다고 말하는 진술은, 관습 규칙의 단순한
경우에서와 같이, 특정한 행위 양식이 실제 실천에서 일반적으로 표준으로
수용되고 있다는 \emphksans{사실(fact)}에 관한 외부 사실 진술이 더 이상 아닐
수 있다. 그것은 이제 수용되었으나 진술되지 않은 승인 규칙을 적용하는
내부 진술일 수 있으며, 대략적으로 말해 ‘그 체계의 유효성 기준들에 비추어
유효하다(valid given the system’s criteria of validity)'는 의미에 불과할
수 있다. 그러나 이러한 점에서, 그리고 다른 점들에서도, 승인 규칙은
체계의 다른 규칙들과는 다르다. 승인 규칙이 존재한다고 단언하는 것은 오직
외부 사실 진술일 수밖에 없다. 종속된 규칙(subordinate rule)은 일반적으로
무시되고 있더라도 유효할 수 있으며, 그 의미에서 ‘존재’할 수 있는 반면,
승인 규칙은 일정한 기준들을 참조하여 법을 식별하는 법원, 공무담당자,
사적 개인들(private persons)의 복합적이지만 통상적으로 일치하는
실천으로서만 존재한다. 그 존재는 하나의 사실 문제이다.

\subsection{2. 새로운 질문들(New
Questions)}\label{uxc0c8uxb85cuxc6b4-uxc9c8uxbb38uxb4e4new-questions}

법체계의 토대들이 법적으로 무제한인 주권자에 대한 복종 습관에 있다고
보는 관점을 버리고, 대신 규칙 체계에 유효성의 기준들을 제공하는 궁극적
승인 규칙이라는 구상(conception)으로 대체하게 되면, 우리는 흥미롭고
중요한 일련의 질문들과 마주하게 된다. 이러한 질문들은 비교적 새로운
것들인데, 그 이유는 오랫동안 법리학과 정치이론이 구시대적인 사유방식에
매여 있었던 동안에는 그러한 질문들이 가려져 있었기 때문이다. 이들은 또한
어려운 질문들이며, 충분한 답변을 위해서는 한편으로는 헌법(constitutional
law)의 몇몇 근본 문제들에 대한 이해가, 다른 한편으로는 법적 형식들(legal
forms)이 조용히 전환되고 변화하는 특유의 방식을 파악할 수 있는 통찰이
요구된다. 따라서 우리는 이 질문들을, 법 개념의 해명에서 일차 규칙들과
이차 규칙들의 결합에 중심적인 위치를 부여해야 한다고 우리가 해왔듯이
고수하는 것이 현명한지 또는 현명하지 않은지와 관련되는 범위 내에서만
탐구할 것이다.

\textbf{(p.~111)} 첫 번째 어려움은 분류(classification)의 문제이다.
왜냐하면 궁극적으로 법을 식별하기(identify) 위해 사용되는 그 규칙은,
흔히 전부를 포괄하는(exhaustive) 것으로 간주되는 법체계를 기술하는
관행적(conventional) 범주들로부터 벗어나 있기 때문이다. 따라서
다이시(Dicey) 이후의 영국 헌법 이론가들은 대체로, 영국의 헌정적 배치들이
한편으로는 엄밀히 그렇게 불리는 법들(laws strictly so called)---즉
제정법들(statutes), 추밀원 명령들(orders in council), 그리고
선례(precedents)에 구현된 규칙들---로, 다른 한편으로는 단지
용례(usages), 양해(understandings), 또는 관습(customs)에 불과한
관행들(conventions)로 구성되어 있다고 반복하여 진술해 왔다. 후자에는,
예컨대 여왕은 상원과 하원에서 적법하게 통과된 법안에 대해 동의를 거부할
수 없다는 것과 같은 중요한 규칙들이 포함된다. 그러나 여왕에게 그러한
동의를 할 법적 책무는 존재하지 않으며, 이러한 규칙들이 관행이라고 불리는
이유는 법원이 그것들을 법적 책무를 부과하는 것으로 인정하지 않기
때문이다. 분명히, 의회 내 여왕이 제정하는 것은 법이다라는 규칙은 이들
어느 범주에도 속하지 않는다. 그것은 법원이 그 규칙과 가장 밀접하게
관련되어 있고 이를 사용하여 법을 식별하는 점에서 관행이 아니며, 동시에
그것이 식별의 기준으로 사용되는 ‘엄밀히 그렇게 불리는 법들’과 동일한
수준의 규칙도 아니다. 설령 이 규칙이 제정법으로 제정되었다고 하더라도,
그것이 단순한 제정법의 수준으로 낮아지는 것은 아니다. 그러한 제정 행위의
법적 지위는 필연적으로, 그 규칙이 해당 제정에 선행하고 그 제정과
독립적으로 존재했다는 사실에 의존하기 때문이다. 더 나아가, 앞 절에서
보였듯이, 그 존재는 제정법의 존재와는 달리 실제 실천(actual
practice)으로 구성되어야 한다.

이러한 측면은 일부에게서 절망의 외침을 끌어낸다: 즉, 분명히 법인 헌법의
근본 규정들이 정말로 법이라는 것을 우리는 어떻게 입증할 수 있는가? 이에
대해 다른 이들은, 법체계의 토대에는 ‘법이 아닌 것’, 즉
‘전법적인(pre-legal) 것’, ‘메타-법적인(meta-legal) 것’, 또는 단순히
‘정치적 사실(political fact)’이 존재한다고 고집한다. 이러한 불안감은
모든 법체계에서 가장 중요한 이 특징을 기술하기 위해 사용되어 온 범주들이
지나치게 거칠다는 확실한 징후이다. 승인 규칙을 ‘법’이라고 부를 수 있다는
근거는, 체계 내의 다른 규칙들을 식별하기 위한 기준을 제공하는
규칙이야말로 법체계를 정의하는 특징(defining feature)으로 충분히 간주될
수 있으며, 따라서 그 자체가 ‘법’이라고 불릴 만하다는 점에 있다. 반대로
그것을 ‘사실’이라고 부를 수 있다는 근거는, 그러한 규칙이 존재한다고
단언하는 것이 기실 \textbf{(p.~112)} ‘실효성 있는(efficacious)’ 체계에서
규칙들이 어떠한 방식으로 식별되는지에 관한 실제 사실(actual fact)에 대한
외부 진술이기 때문이다. 이 두 측면은 모두 주목을 요구하지만, 우리는 ‘법’
또는 ‘사실’이라는 어느 한 명칭(label)을 선택하는 것만으로는 이 둘 모두를
정당하게 다룰 수 없다(do justice). 대신, 우리는 궁극적 승인 규칙을 두
가지 관점에서 바라볼 수 있음을 기억해야 한다. 하나는 그 규칙이 체계의
실제 실천(actual practice) 속에 존재한다는 외부 사실 진술로 표현되는
관점이며, 다른 하나는 그 규칙을 사용하여 법을 식별하는 사람들이 행하는
유효성(validity)에 관한 내부 진술로 표현되는 관점이다.

두 번째 질문군은 특정 국가나 사회 집단 안에 법체계가
\emphksans{존재한다(exists)}고 하는 단언(assertion)의 숨겨진 복잡성과
모호성(vagueness)에서 비롯된다. 우리가 이러한 단언을 할 때, 실제로는
대개 함께 수반되는 서로 이질적인 여러 사회적 사실들(heterogeneous social
facts)을 압축적 포괄 형식(portmanteau form)으로 가리키고 있는 것이다.
오해를 불러일으키는 이론의 그림자 아래 발전된 법적 사고와 정치적 사고의
표준 용어법(standard terminology)은 사실들을 지나치게 단순화하고 가리기
쉽다. 그러나 이러한 용어법이 구성하는 안경(spectacles)을 벗고 사실들을
바라보면, 법체계(legal system)는 인간과 마찬가지로, 어떤 단계에서는 아직
태어나지 않은 상태(unborn)에 있고, 또 어떤 단계에서는 아직 모체로부터
완전히 독립하지 않은 상태(not yet wholly independent)에 있으며, 이후
건강하고 독립적인 존재(healthy independent existence)를 누리다가,
쇠퇴(decay)하고 마침내 사망(die)하기도 한다는 것이 분명해진다. 이러한
‘출생과 정상적 독립 존재 사이’, 또는 ‘그 독립된 존재와 죽음 사이’의 중간
단계들은 우리가 익숙하게 사용하는 법 현상에 대한 설명 방식을 어긋나게
만든다(put out of joint). 그러나 이 중간 단계들은 그만큼 우리가 정상
사례에서 “특정 국가에는 법체계가 존재한다”라는 확신에 찬 참된 단언을
할 때 당연하게 여기는 것의 온전한 복잡성을 부각시키기 때문에,
혼란스러울지라도 연구할 가치가 있다.

이러한 복합성을 인식하는 한 가지 방식은, 명령들에 대한 일반적 복종의
습관이라는 단순한 오스틴식 공식(Austinian formula)이, 사회가 법체계를
갖기 위해 충족해야 하는 최소 조건들을 구성하는 복합적 사실들을 재현하는
데서 정확히 어디에서 실패하거나 왜곡하는지를 살펴보는 데 있다. 우리는 이
공식이 하나의 필요조건을 지목한다(designate)는 점은 인정할 수 있다. 즉,
법들이 의무(obligation)나 책무(duty)를 부과하는 경우, 그 법들에는
일반적으로 복종이 이루어져야 하거나 적어도 일반적으로 불복종이
이루어지지 않아야 한다는 점이다. 그러나 이는 필수적이긴 하나, 법체계가
사적 시민에게 영향을 미치는 지점에서 나타나는 우리가 ‘최종 산물(end
product)’이라 부를 수 있는 것만을 다룰 뿐이며, 법체계의 일상적 존재는
\textbf{(p.~113)} 또한 공무담당자들에 의한 법의 창출(official creation),
공무담당자들에 의한 식별, 그리고 공무담당자들에 의한 사용 및 적용으로도
이루어져 있다. 여기에서 관련되는 법과의 관계를 ‘복종(obedience)’이라
부를 수 있으려면, 그 단어를 통상적 용법을 훨씬 넘어 확장하여, 결국에는
이러한 작용들을 정보를 주는 방식으로 특징짓는 기능을 상실하게 해야 한다.
입법자들이 법들을 제정할 때 자신들의 입법 권한들을 부여하는 규칙들에
부합하는 경우, 그러한 권한들을 부여하는 규칙들이 그것들을 따를 책무를
부과하는 규칙들에 의해 강화되지 않는 한, 어떤 통상적 의미에서도
입법자들이 규칙에 ‘복종’한다고 말할 수는 없다. 또한 이러한 규칙들에
부합하지 않았다고 해서 그들이 법에 ‘불복종하는(disobey)’ 것도 아니며,
다만 법을 만들어내는 데 실패할 뿐이다. 마찬가지로, 판사가 체계의 승인
규칙(rule of recognition)을 적용하여 어떤 제정법(statute)을 유효한
법으로 인정하고 분쟁의 판정(determination)에 이를 사용하는 경우에,
‘복종’이라는 말은 그가 하는 일을 적절히 묘사하지 못한다. 물론 우리는
사실들 앞에서 ‘복종’이라는 단순한 용어를 유지하고자 한다면, 여러 장치를
통해 그렇게 할 수는 있다. 그 한 예로, 판사가 제정법을 인정하는 데서
사용하는 일반적 유효성 기준을, ‘헌법 제정자들(Founders of the
Constitution)’이 내린 명령들에 대한 복종의 사례로 표현하거나, (그러한
‘제정자들’이 없는 경우에는) 지휘관 없는 사령, 즉 ‘탈심리화된
사령(depsychologized command)’에 대한 복종으로 표현할 수 있을 것이다.
그러나 후자의 관념은, 삼촌 없는 조카라는 관념만큼이나, 우리의 주의를
진지하게 요구할 이유가 없을지도 모른다. 다른 한편으로는, 법의 공무담당
측면 전체를 시야에서 밀어내고 입법과 재판에서 이루어지는 규칙의 사용에
대한 서술을 포기한 채, 공무담당 세계 전체를 하나의 인격(‘주권자’)이 여러
대리인이나 대변자를 통해 명령들을 발하는 것으로 생각하고, 시민이 그
명령들에 습관적으로 복종한다고 간주할 수도 있다. 그러나 이는 여전히
서술을 기다리는 복합적 사실들에 대한 편의적 축약에 불과하거나, 아니면
재앙적으로 혼란을 초래하는 신화적 구성물일 뿐이다.

법체계가 존재한다는 것을, 기실 일반 시민과 법의 관계를 특징짓기는 하지만
그것을 언제나 남김없이 서술하지는 않는 습관적 복종(habitual
obedience)이라는 기분 좋을 만큼 단순한 표현으로 설명하려는 시도가 실패한
데 대한 반작용으로, 그 반대의 오류를 범하는 일은 자연스러운 일이다. 그
오류는 공무담당 활동들(official activities)---특히 사법적 태도(judicial
attitude)나 법에 대한 관계---의 특징이지만 역시 남김없는 설명은 아닌
것을 취하여, 그것을 법체계를 가진 사회집단에서 반드시 존재해야 하는 것에
대한 충분한 설명으로 취급하는 데 있다. 이러한 오류는 사회 구성원의
대다수가 법에 대해 단순히 습관적으로 복종한다는 단순한 구상을, 그들이
법들의 유효성을 최종적으로 평가하는 기준을 특정하는 궁극적 승인
규칙(rule of recognition)을 일반적으로 공유하거나, 수용하거나, 구속력
있는 것으로 간주해야 한다는 구상으로 대체하는 데 이른다. 물론 우리는
3장에서 살펴본 것처럼, 법의 원천들에 관한 지식과 이해가 사회 전반에 널리
퍼진 단순한 사회를 상상할 수 있다. 그곳에서는 ‘헌법(constitution)’이
너무 단순해서, 일반 시민뿐 아니라 공무담당자들과 법률가들 역시 그 헌법을
알고 있으며 수용하고 있다고 말해도 아무런 허위가 없을 것이다. 렉스 1세의
단순한 세계에서는, 인구 다수가 그의 말에 단순히 습관적으로 복종한 것 그
이상이 있었다고 말할 수도 있다. 그 사회에서는 수범자들과 체계의
공무담당자들 모두가 렉스의 말(Rex's word)을 사회 전체를 위한 유효한 법의
기준으로 규정하는 하나의 승인 규칙을 명시적이고 의식적인 방식으로
‘수용하고(accepted)’ 있었을 수 있다. 물론 수범자들과 공무담당자들은
수행해야 할 역할과 이 기준에 따라 식별되는 법규칙들에 대한 관계가 서로
달랐겠지만 말이다. 그러나 이렇게 단순한 사회에서 가능한 상태를, 복잡한
현대 국가에서도 항상 또는 일반적으로 존재한다고 주장하는 것은 허구를
고집하는 일일 것이다. 여기서 상황의 현실은 분명 이렇다: 상당수의 일반
시민, 어쩌면 과반수는 법적 구조나 그 유효성 기준들에 대한 일반적 구상을
가지고 있지 않다. 그가 복종하는 당해 법(the law)은 그에게 단지
‘법’으로만 알려져 있는 어떤 것일 뿐이다. 그는 여러 서로 다른 이유들로 그
법에 복종할 수 있으며, 그 이유들 중에는 그러는 것이 자신에게 가장
이롭다는 지식이 자주, 비록 항상은 아니지만, 포함될 수 있다. 그는 불복종
시 발생할 수 있는 일반적으로 개연성 있는 결과들---예컨대 자신을 체포할
수 있는 공무담당자들이 있고, 자신을 재판하여 법을 어겼다는 이유로 감옥에
보낼 수 있는 이들이 있다는 것---을 알고 있을 것이다. 그 체계의 유효성
검사 기준들에 따라 유효한 법들에 인구 대다수가 복종하는 한, 이는 특정한
법체계가 존재한다는 것을 확립하기 위해 필요한 증거 전부이다.

그러나 법체계가 일차 규칙들과 이차 규칙들의 복합적 결합이라는 바로 그
이유 때문에, 이러한 증거만으로는 법체계의 존재에 수반되는 법과의
관계들을 모두 서술하기에 충분하지 않다. 이는 체계의 공무담당자들이
공무담당자로서 관련되는 이차 규칙들에 대해 갖는 관계에 관한
\textbf{(p.~115)} 서술로 보충되어야 한다. 여기서 핵심적인 것은 체계의
유효성 기준들을 포함하는 승인 규칙에 대한 공무담당자들의 통일적이거나
공유된 수용(a unified or shared official acceptance)이 존재해야 한다는
점이다. 그러나 바로 이 지점에서, 일반 시민의 경우에 불가결한 최소를
특징짓기에는 충분했던 일반적 복종이라는 단순한 관념은 부적절해진다.
요점은 ‘복종(obedience)’이라는 표현이 법원과 다른 공무담당자들이 이러한
이차 규칙들을 규칙으로서 존중하는 방식을 가리키는 데 자연스럽게 사용되지
않는다는 ‘언어적(linguistic)’ 문제가 아니며, 적어도 단지 그 문제만은
아니다. 필요하다면, 우리는 ‘따르다(follow)’, ‘준수하다(comply)’, 또는
‘부합하다(conform to)’와 같은 보다 포괄적인 표현을 찾아낼 수 있을
것이다. 이러한 표현들은 시민이 병역에 응하기 위해 법과 관련하여 수행하는
행위와, 판사가 의회 내 여왕이 제정하는 것은 법이라는 전제하에 자신의
법정에서 특정 제정법(statute)을 법으로 식별하는 행위를 모두 특징지을 수
있을 것이다. 그러나 이러한 포괄 용어들(blanket terms)은, 우리가
법체계라고 부르는 이 복합적 사회 현상의 존재에 포함된 최소 조건을
이해하기 위해 반드시 파악되어야 할 중대한 차이들을 단지 가려버릴 뿐이다.

입법자들이 자신들의 권한들을 부여하는 규칙들에 부합하여 입법 활동을
하거나, 법원이 수용된 궁극적 승인 규칙(ultimate rule of recognition)을
적용할 때, 이를 ‘복종(obedience)’이라고 묘사하는 것이 오해를
불러일으키는 이유는 다음과 같다. 어떤 규칙(또는 명령)에 복종한다고 해서,
그 복종하는 사람이 자신이 하는 일이 자기 자신뿐 아니라 다른 사람들에게도
옳은 일이라고 생각해야 할 \emphksans{필요는(need)} 없다. 그는 자신이 하는
행동을 사회집단의 다른 구성원들에게도 적용되는 행동 기준을 실현하는
것으로 보지 않아도 된다. 그는 자신이 보이는 부합 행위를 ‘옳다(right)’,
‘올바르다(correct)’, 또는 ‘의무적이다(obligatory)’라고 생각할 필요가
없다. 즉, 그의 태도에는 사회 규칙이 수용되거나 특정 행위 유형이 일반적
기준으로 간주될 때 수반되는 비판적 성격(critical character)이 전혀 없을
수 있다. 그는 그러한 규칙들을 그 규칙들이 적용되는 모든 사람을 위한
기준으로 수용하는 내부 관점(internal point of view)을 공유하지 않을 수도
있다. 대신 그는 그 규칙을 단지 \emphksans{그에게(him)} 요구되는 것, 즉
제재적-불이익(penalty)의 위협 아래 행동을 요구하는 것으로만 여길 수
있다. 그는 그저 결과에 대한 두려움 때문에, 혹은 관성에 따라 복종할 수
있으며, 자신이나 타인이 그것을 따를 의무가 있다고 여기지 않고, 규칙을
어기는 자신이나 타인을 비판할 의향도 없을 수 있다. 하지만,
\textbf{(p.~116)} 모든 일반 시민이 규칙에 복종할 때 가질 \emphksans{수도
있는(may)} 이러한 개인적인 차원의 관심만으로는, 법원이 자신이 법원으로서
다루는 규칙들에 대해 갖는 태도를 설명할 수 없다. 이는 특히 다른 규칙들의
유효성을 평가하는 데 쓰이는 궁극적 승인 규칙(ultimate rule of
recognition)의 경우에 분명하다. 이 규칙은, 존재하기 위해서 반드시 하나의
공적인(public), 공통의(common) 올바른 사법적 결정의 기준(standard of
correct judicial decision)으로서 내부 관점(internal point of view)에서
수용되어야 하며, 단순히 각 판사가 자기 몫으로만 ‘복종’하는 어떤 것으로
여겨져서는 안 된다. 그 법체계의 법원들 중 개별적인 법원이 때때로 이
규칙들로부터 이탈하는 일이 있을 수는 있으나, 일반적으로 그러한 이탈을
공통적이거나 공적인 기준(common or public standards)에서의
일탈(lapses)로 비판적으로 다루어야 한다. 이는 단지 법체계의 효율성이나
건전성과 관련된 문제가 아니라, 하나의 법체계(single legal system)가
존재한다고 말할 수 있기 위한 논리적으로 필수적인 조건(logically
necessary condition)이다. 가령 어떤 판사들만이 ‘자기 몫으로만’ 행동하며
“의회 내 여왕이 제정하는 것은 법이다”라는 전제에 서고, 이 승인 규칙을
존중하지 않는 이들을 아무도 비판하지 않는다면, 법체계의 특유의 통일성 및
계속성(characteristic unity and continuity)은 사라지고 말 것이다. 이러한
통일성과 계속성은 바로 이 핵심적 지점에서 법적 유효성에 관한 공통된
기준(common standards of legal validity)이 수용되는 데 의존하기
때문이다. 이러한 사법적 행태의 변칙들과, 평범한 사람이 서로 상반되는
사법 명령들에 직면하게 될 때 결국 초래될 혼란 사이의 그 사이 구간에서,
우리는 그 상황을 어떻게 묘사해야 할지 난처할 것이다. 우리는 오직 그것이
평소에는 너무 자명해서 인식되지 않는 것들을 날카롭게 의식하게 해주기
때문에만 생각해볼 만한 가치가 있는 하나의 \emphksans{자연의 기이함(lusus
naturae)}과 마주하게 될 것이다.

따라서 하나의 법체계가 존재하기 위해 필요하고 충분한 최소 조건은 두
가지이다. 첫째, 그 법체계의 궁극적 유효성 기준들(ultimate criteria of
validity)에 따라 유효한 것으로 간주되는 행위 규칙들(rules of
behaviour)에는 일반적으로 복종이 이루어져야 한다. 둘째, 법적 유효성
기준들을 특정하는 승인 규칙들(rules of recognition)과 변경 규칙들(rules
of change) 및 재판 규칙들(rules of adjudication)은, 해당 체계의
공무담당자들 사이에서 공무담당 행위의 공통된 공적 기준들(common public
standards of official behaviour)로 실효적으로 수용되어야 한다. 첫 번째
조건은 일반 시민들이 만족시킬 \emphksans{필요가 있는(need)} 유일한 조건이다:
그들은 각자 ‘자기 몫으로만’ 복종할 수 있으며, 그 동기는 어떠하든
상관없다; 물론 건강한 사회에서는 실제로 많은 시민들이 이 규칙들을 공통의
행동 기준들로 수용하고, 그것들에 복종할 의무가 있음을
인정하거나(acknowledge), 심지어 이러한 의무를 \textbf{(p.~117)} 헌법을
존중할 보다 일반적인 의무로 거슬러 추적하기도 할 것이다. 두 번째 조건은
법체계의 공무담당자들이 반드시 만족시켜야 한다. 그들은 해당 규칙들을
공무담당 행위의 공통된 기준들(common standards of official behaviour)로
간주하고, 자기 자신과 동료들의 일탈(deviations)을 기준에서의
이탈(lapses)로 비판적으로 평가해야 한다. 물론 여기에 덧붙여,
공무담당자가 단지 개인적 행위능력(personal capacity)에서 복종하기만 하면
되는 일차 규칙들도 많다는 점은 사실이다.

따라서 “법체계가 존재한다”는 단언은 이중적인 얼굴을 가진(Janus-faced)
진술이다. 한편으로는 일반 시민들의 복종을, 다른 한편으로는
공무담당자들이 이차 규칙들을 공무담당 행위의 비판적 공통 기준들(critical
common standards of official behaviour)으로 수용하는 것을 함께 가리킨다.
이러한 이중성에 놀랄 필요는 없다. 이것은 단지 법체계의 복합적
성격(composite character)---즉 단순하고 분권적인 전법적 사회 구조와
대비되는---을 반영하는 것일 뿐이다. 단지 일차 규칙들만으로 구성된 더
단순한 사회 구조에서는, 공무담당자가 존재하지 않기에, 해당 규칙들은
집단의 행위에 대한 비판적 기준으로서 널리 수용되어야 한다. 만약 그러한
경우에 내부 관점(internal point of view)이 널리 퍼져 있지 않다면,
논리적으로 규칙이 존재할 수 없다. 하지만, 우리가 주장해온 바와 같이,
법체계를 이해하는 가장 유익한 방식은 일차 규칙들과 이차 규칙들의
결합(union of primary and secondary rules)으로 보는 것이며, 이 경우 집단
전체의 공통 기준으로서 규칙을 수용하는 일은, 일반 개인이 자신 몫으로
복종함으로써 그 규칙에 묵인하는 비교적 수동적인 태도와 분리될 수 있다.
극단적인 경우에는, 법적 언어의 특유의 규범적 사용(예: “이것은 유효한
규칙이다”)을 포함한 내부 관점이 오직 공무담당 세계 내부에만 존재할 수도
있다. 이처럼 보다 복잡한 체계에서는, 오직 공무담당자들만이 법적 유효성에
대한 체계의 기준들을 수용하고 사용할 수 있다. 그러한 사회는 비참할 만큼
양처럼 순응적인(sheeplike) 사회일 수 있고, 결국 도살장으로 향하게 될지도
모른다. 하지만 그 사회가 실제로 존재할 수 없다고 보거나, 그것이
법체계라는 명칭을 가질 수 없다고 단정할 이유는 별로 없다.

\subsection{3. 법체계의 병리학(THE PATHOLOGY OF A LEGAL
SYSTEM)}\label{uxbc95uxccb4uxacc4uxc758-uxbcd1uxb9acuxd559the-pathology-of-a-legal-system}

따라서 법체계의 존재에 대한 증거는 사회생활의 두 가지 서로 다른
영역으로부터 도출되어야 한다. 우리가 법체계가 존재한다고 확신을 가지고
말할 수 있는 정상적이고 문제없는 경우란, 바로 이 두 영역이
\textbf{(p.~118)} 법에 대한 각자의 전형적인 관심사에 있어서 서로
일치(congruent)하고 있음이 분명한 경우이다. 거칠게 말하면, 사실관계는
다음과 같다. 즉, 공무담당자 수준에서 유효한 것으로 인정되는 규칙들에
일반적으로 복종이 이루어진다. 그러나 때로는 공무담당 영역이 사적
영역으로부터 분리될(detached) 수 있는데, 이는 법원에서 사용되고 있는
유효성 기준들(criteria of validity)에 따라 유효한 규칙들에 더 이상
일반적 복종이 존재하지 않는다는 의미에서 그러하다. 이러한 상황이
발생하는 다양한 방식들은 법체계의 병리(pathology of legal systems)에
속한다. 왜냐하면 그것들은, 우리가 법체계가 존재한다는 외부 사실
진술(external statement of fact)을 할 때 지칭되는(referred to)
복합적이고 상호 일치적인 실천의 붕괴를 나타내기 때문이다. 여기에는,
특정한 체계 내부로부터 법에 관한 내부 진술을 할 때마다
전제(presupposed)되는 것의 부분적 실패가 존재한다. 이러한 붕괴는 서로
다른 교란 요인들의 산물일 수 있다. 집단 내부에서 통치에 대한 경쟁적
주장들이 제기되는 ‘혁명(revolution)’은 그 한 사례에 불과하며, 이는
언제나 기존 체계의 일부 법들의 위반을 수반하겠지만, 새로운 헌법이나
새로운 법체계가 아니라 단지 공무담당자들을 새로운 개인 집합으로 법적
권한 없이 대체하는 것만을 수반할 수도 있다. 기존 체계 하에서의 권한 없이
외부로부터 통치에 대한 경쟁적 주장이 제기되는 적군 점령(enemy
occupation)도 또 하나의 사례이며, 통치에 대한 정치적 주장조차 없이
무정부 상태나 도적 행위에 직면하여 질서 있는 법적 통제가 단순히 붕괴되는
경우 역시 또 다른 사례이다.

이러한 각각의 경우에서, 법원이 해당 영토에서든 망명 상태에서든
기능하면서, 한때 안정적으로 확립된 옛 체계의 유효성 기준을 여전히
사용하는 중간 단계들(half-way stages)이 있을 수 있다. 그러나 이러한
법질서들은 해당 영토에서 실효적이지 않다. 이러한 경우에 “법체계가
완전히 소멸되었다”고 말할 수 있는 시점은, 어떤 정확한 확정(exact
determination)도 허용하지 않는 사안이다. 명백하게, 복구될 상당한
가망(chance)이 있거나, 기존 체계의 혼란이 결과가 아직 확정되지 않은
전면적 전쟁의 일부 사건일 경우에는, 그 체계가 더 이상 존재하지 않는다는
무조건적 단언(assertion)은 정당화될 수 없다. 이는 바로 “법체계가
존재한다”는 진술(statement)이 중단을 허용할 만큼 충분히 폭넓고 일반적인
유형이기 때문이다. 이 진술은 짧은 시간 동안 일어나는 일들에 의해
검증되거나 반증되는 것이 아니다.

물론 이러한 중단 이후에 법원과 인구 간의 정상적인 \textbf{(p.~119)}
관계가 회복된 뒤에는 어려운 질문들이 제기될 수 있다. 점령군의 축출이나
반란 정부의 패배 이후 망명 정부가 귀환하면, 그 중단의 시기에 해당
영토에서 무엇이 ‘법’이었는지, 혹은 무엇이 아니었는지에 대한 질문들이
생긴다. 여기서 가장 중요한 점은, 이 질문이 사실(fact)의 문제가 아닐 수도
있다는 점이다. 만약 그것이 사실의 문제라면, 해당 중단이 너무나
장기적이고 완전했기 때문에 기존 법체계가 소멸되었고 망명 귀환 시 유사한
새로운 체계가 수립되었다고 기술해야 하는지를 묻는 방식으로
확정되어야(settled) 할 것이다. 그러나 실제로는 이 질문은 국제법상의
문제로 제기되거나, 다소 역설적으로, 복귀 이후 존재하게 된 바로 그 법체계
내부의 법 문제로 제기될 수 있다. 후자의 경우, 복귀한 체계가, 그 체계가
계속해서 해당 영토의 법이었다고 소급적으로 선언하는 법을 포함할 수도
있다---혹은 좀 더 솔직히 말해, 그 체계가 해당 시기 동안 법이었던 것으로
간주된다(deemed)고 선언할 수 있다. 이러한 선언은, 그 중단 기간이 사실
문제로 취급되었을 경우 도달했을 결론과는 상당히 불일치해 보일 수
있음에도 불구하고, 그렇게 이루어질 수 있다. 그런 경우에, 해당 선언은
중단 기간 중 발생한 사건들과 거래들에 대해 법원이 적용해야 할 법을
정하는 복귀 체계의 하나의 규칙으로서, 얼마든지 존립할(stand) 수 있다.

여기에 역설이 있는 것은 오직, 법체계가 스스로의 과거, 현재, 또는 미래의
존재 단계로 간주되어야 하는 것에 관하여 내리는 법 진술들(statements of
law)을, 외부 관점에서 제시되는 그 존재에 대한 사실 진술과
경쟁하는(rivals) 것으로 생각할 때뿐이다. 자기-참조(self-reference)의
겉보기 수수께끼를 제외하면, 어떤 현행 체계 안에서 그 체계가 존재한
것으로 간주되어야 하는 기간에 관한 규정의 법적 지위는, 한 체계의 법이
다른 나라에 어떤 체계가 아직 존재한다고 선언하는 경우와 다르지 않다.
물론 후자의 경우 실천적 귀결이 많을 개연성은 높지 않다. 예를 들어,
우리는 소련 영토에서 현재 존재하는 법체계가 사실상 차르 체제의 법체계가
아니라는 점을 분명히 알고 있다. 그러나 만약 영국 의회의
제정법(statute)이 차르 시대 러시아의 법이 여전히 러시아 영토의 법이라고
선언한다면, 이 제정법은 소련(USSR)에 관한 영국법의 일부로서 기실 의미와
법적 효과를 가질 것이다. \textbf{(p.~120)} 그러나 이는 직전 문장에 담긴
사실 진술의 참임에 전혀 영향을 미치지 않는다. 그 제정법의 효력(force)과
의미는, 영국 법원이 적용할 법을 규정하고, 그 결과로 영국 내에서 러시아와
관련된 사안에 적용될 법을 정하는 것에 있을 뿐이다.

방금 서술한 상황의 역상황(converse)은, 새로운 법체계가 옛 법체계의
자궁으로부터---때로는 제왕절개(Caesarian operation)를 거친
뒤에야---출현하는 매혹적인 전환기의 순간들에서 볼 수 있다.
영연방(Commonwealth)의 최근 역사는 법체계 발생론(embryology of legal
systems)의 이러한 측면을 연구하기에 훌륭한 장(場)을 제공한다. 이 발전
과정을 도식적이고 단순화하여 개관하면 다음과 같다. 한 시기의 초기에
우리는 지방 입법부(local legislature), 사법부(judiciary),
행정부(executive)를 갖춘 식민지를 발견할 수 있다. 이러한 헌정 구조는
영국 의회(United Kingdom Parliament)의 제정법에 의해 설정된 것으로, 영국
의회는 식민지에 대해 입법할 완전한 법적 권능(legal competence)을
보유하고 있다. 여기에는 지방 법들(local laws)뿐 아니라, 식민지 헌법에
관한(referring to) 제정법들을 포함하여, 자기 자신의 제정법들(statutes)을
개정하거나 폐지할 권한도 포함된다. 이 단계에서 식민지의 법체계는, 의회
내 여왕(Queen in Parliament)이 제정하는(enact) 것은 \emphksans{(다른 것들
중에서도(inter alia))} 식민지에 대하여 법이다라는 궁극적 승인 규칙(rule
of recognition)에 의해 특징지어지는 더 넓은 체계의 명백한
종속된(subordinate) 부분이다. 발전 단계의 말미에 이르면, 궁극적 승인
규칙이 전환되었음(shift)을 확인하게 되는데, 이는 웨스트민스터
의회(Westminster Parliament)가 옛 식민지에 대해 입법할 법적 권능(legal
competence)이 더 이상 그곳의 법원들에서 인정되지 않기 때문이다. 옛
식민지의 헌정 구조 상당 부분이 여전히 웨스트민스터 의회의 원래
제정법(statute)에서 발견된다는 점은 여전히 사실이다. 그러나 이는 이제
단지 역사적 사실에 불과하며, 그것은 더 이상 해당 영토에서의 현재적 법적
지위를 웨스트민스터 의회의 권위에 빚지고 있지 않다. 옛 식민지의 법체계는
이제 ‘지역적 뿌리(local root)’를 가지게 되는데, 이는 법적 유효성의
궁극적 기준들을 특정하는 승인 규칙이 더 이상 다른 영토의 입법부가 제정한
것들(enactments)을 참조(refer)하지 않기 때문이다. 새로운 규칙은, 그것이
해당 지역 체계의 사법적 및 기타 공무담당 작용에서 그러한 규칙으로
수용되고 사용된다는 사실에 단순히 기초해 있다. 이 지역 체계의 규칙들에는
일반적으로 복종이 이루어진다. 따라서 비록 지방 입법부의 구성, 제정 방식,
구조가 여전히 원래 헌법에서 규정된 바와 동일할 수 있다 하더라도, 그
제정물들(enactments)이 현재 유효한 것은 \textbf{(p.~121)} 더 이상
웨스트민스터 의회의 유효한 제정법(statute)이 부여한 권한들의 행사이기
때문이 아니다. 그것들은, 지역적으로 수용된 승인 규칙에 따라, 지방
입법부의 제정이 유효성의 궁극적 기준이기 때문에 유효한 것이다.

이러한 발전은 여러 방식으로 달성될 수 있다. 모국 입법부는 일정 기간 동안
사실상 식민지의 동의 없이는 식민지에 대한 자신의 형식적 입법 권위(formal
legislative authority)를 행사하지 않다가, 마침내 해당 옛 식민지에 대한
입법 권한을 포기함으로써 무대에서 퇴장할 수 있다. 이때 주목할 점은,
웨스트민스터 의회가 이와 같이 스스로의 권한들을 불가역적으로 줄일(cut
down) 법적 권능(legal competence)을 가지고 있는지를 영국 법원이
인정할지에 대해 이론적인 의문이 존재한다는 것이다. 반대로, 이러한 단절은
폭력을 통해서만 이루어질 수도 있다. 그러나 어느 경우이든, 이 발전의
말미에는 두 개의 독립된 법체계가 존재하게 된다. 이것은 하나의 사실
진술(factual statement)이며, 비록 그것이 법체계의 존재에 관한 것이라
할지라도, 그렇다고 해서 덜 사실적인 것은 아니다. 이에 대한 주된 증거는,
이전 식민지에서는 현재 수용되고 사용되는 궁극적 승인 규칙이 더 이상 타
지역의 입법부 작용에 대한 어떠한 참조도 유효성 기준들 가운데 포함하지
않는다는 사실이다.

그런데 다시, 영연방(Commonwealth)의 역사에서 흥미로운 사례들이
보여주듯이, 사실상 식민지의 법체계가 이제 모국으로부터 독립되었음에도
불구하고, 모국의 법체계는 이 사실을 인정하지 않을 수도 있다.
웨스트민스터 의회가 식민지에 대해 입법할 권한을 보유하고 있거나 법적으로
되찾을 수 있다는 것이 여전히 영국법의 일부일 수 있으며, 실제로 그러한
경우에, 만약 웨스트민스터 제정법(Westminster statute)과 현지 입법부의 법
하나가 충돌하는 사건이 영국 국내 법원에 회부된다면, 영국 법원은 이
사태에 대한 이러한 관점에 효력을 부여할 수도 있다(give effect). 이 경우,
영국법명제들(propositions of English law)은 사실과 충돌하는 듯 보인다.
식민지의 법은 사실상 독립된 법체계이며, 자체의 지역적 궁극적 승인 규칙을
가진 체계인데, 영국 법원은 이를 사실상 그러한 바로 그것으로 인정하지
\emphksans{않는다(not)}. 사실상은 두 개의 법체계가 존재하지만, 영국법은 오직
하나만 존재한다고 고집한다(insist). 그러나, 한쪽 단언은 사실
진술(statement of fact)이고 다른 쪽은 (영국)법명제(proposition of
law)이기 때문에, 이 둘은 논리적으로 충돌하지 않는다. 이 입장을 명확히
하기 위해, 우리는 사실 진술이 참이며(true), 영국법명제는 “영국법상
올바른(correct in English law)” 것이라고 말할 수 있다.
\textbf{(p.~122)} 이러한 독립된 법체계의 존재 여부에 대한 사실적
단언(또는 부정)과 법체계의 존재에 관한 법명제들 간의 구별은, 국제공법과
국내법 사이의 관계를 논할 때에도 반드시 유념해야 한다. 어떤 매우 기이한
이론들이 그럴듯해 보이는 유일한 이유는 이 구별을 무시하기 때문이다.

이러한 법체계의 병리학(pathology)과 발생학(embryology)에 대한 거친
개관을 마무리하려면, 이제 우리는, “하나의 법체계가 존재한다”는
무조건적 단언(assertion)에 의해 그 일치(congruence)가 단언되는 정상
조건들이 부분적으로 실패하는 다른 형태들도 주목해야 한다. 법체계
내부에서 법에 대한 내부 진술이 행해질 때 일반적으로 전제되는,
공무담당자들 사이의 통일성(unity)은 부분적으로 붕괴될 수 있다. 헌법적
문제들---그리고 오직 그 문제들에 한해서---에 대해, 공무담당 세계
내부에서 분열이 발생하고 궁극적으로 사법부 내부의 분열로 이어질 수 있다.
법을 식별하는 데 사용되는 궁극적 기준들에 대한 이러한 분열의 시작은
1954년 남아프리카 공화국에서 발생한 헌법적 위기에서 볼 수 있었으며, 이는
\emph{Harris} 대 \emph{Dönges} 사건으로 법원에 회부되었다. 여기서
입법부는 자신이 보유한 법적 권능(legal competence)과 권한들(powers)에
대해 법원과는 다른 견해를 취하며 입법 조치들을 제정하였고, 법원은 그
조치들을 유효하지 않다고(invalid) 선언하였다. 이에 대한 대응으로
입법부는, 입법부의 제정물들(enactments)을 유효하지 않게 한 통상 법원들의
판결들에 대해 항소가 제기될 수 있도록 특별 항소 ‘법원’을 창설하였다. 이
법원은 실제로 그러한 항소를 심리하고 통상 법원들의 판결들을 번복했으며,
다시 통상 법원들은 그 특별 법원들을 창설한 입법을 유효하지 않다고
선언하고 해당 판결들을 법적으로 무효(nullity)라고 판시하였다. 이 절차가
중단되지 않았더라면(정부가 자신의 뜻을 관철하는 이 수단을 계속 추구하는
것이 지혜롭지 못하다고 판단했기 때문에 중단되었다), 우리는 입법부의
권능(competence)과 그에 따라 유효한 법의 기준들에 대한 두 견해 사이에서
무한 반복되는 진동 현상을 보았을 것이다. 법체계의 승인 규칙을 식별하는
것이 오직 그 조건 아래에서만 가능한 공무담당자들 사이의---특히
사법적---조화라는 정상 조건은 일시적으로 정지되었을 것이다. 그럼에도
불구하고, 이 헌법적 문제와 직접적으로 관련되지 않은 대부분의 법적
작용들은 이전과 마찬가지로 계속되었을 것이다. 인구가 분열되고, ‘법과
질서’가 \textbf{(p.~123)} 붕괴되기 전까지는, 원래의 법체계가 소멸했다고
말하는 것은 오해를 불러올 수 있다. 왜냐하면 ‘동일한 법체계(same legal
system)’라는 표현은 너무 넓고 유연하기 때문에, 유효한 법의
\emphksans{모든(all)} 원래 기준들에 대한 공무담당자들의 통일적 합의가
유지되는 것을 그 법체계가 ‘동일한 것’으로 남아 있기 위한 필요조건으로
삼을 수는 없기 때문이다. 우리가 할 수 있는 일은 지금까지 기술해온
방식대로 해당 상황을 설명하고, 그것을 하나의 표준 이하(substandard),
비정상적인 사례로 기록해 두는 것이다. 이 사례는 그 속에 법체계가
해체될(dissolve) 수 있는 잠재적 위협을 품고 있다.

이 마지막 사례는, 우리가 다음 장에서 논의하게 될 보다 넓은 주제의 경계로
우리를 이끈다. 이는 법체계의 궁극적 유효성 기준이라는 중대한 헌법적
사안(high constitutional matter)과, 그 법체계의 ‘일상적(ordinary)’ 법
양자 모두와 관련된다. 모든 규칙에는 특정 사례들을 일반적 용어의 사례로
식별하거나 분류하는 과정이 포함되며, 우리가 어떤 것을 규칙(rule)이라고
부를 준비가 되어 있는 모든 경우에 있어서, 해당 규칙이 분명히 적용되는
명확한 중심 사례들과, 그 적용 여부를 단언할 이유와 부정할 이유가 모두
존재하는 경계 사례들을 구별하는 것이 가능하다. 우리가 구체적인 상황들을
일반 규칙 아래로 포섭할 때, 이와 같은 확실성의 핵심(core of certainty)과
의문의 반음영(penumbra of doubt)이라는 이중성은 결코 제거될 수 없다.
이로 인해 모든 규칙에는 모호함의 가장자리(fringe), 곧 ‘개방적 직조(open
texture)’가 부여되며, 이는 법을 식별하는 데 사용되는 궁극적 기준들을
명시하는 승인 규칙(rule of recognition)에도, 하나의 특정한
제정법(statute)에도 마찬가지로 영향을 미칠 수 있다. 법의 이러한 측면은
흔히, 규칙이라는 틀로 법 개념을 해명하려는 시도가 오해를 낳을 수 있다는
점을 보여주는 것으로 여겨진다. 상황의 현실에도 불구하고 이를 고수하는
것은 흔히 ‘개념주의(conceptualism)’나 ‘형식주의(formalism)’라는 낙인이
찍히곤 한다. 우리는 이제 이와 같은 비판을 평가하는 논의로 넘어가게 될
것이다.

\newpage

\subsection{제6장 주석}\label{uxc81c6uxc7a5-uxc8fcuxc11d}


{\footnotesize
\noteentry 100쪽. \notetitle{승인 규칙(rule of recognition)과 켈젠(Kelsen)의 ‘기초규범(basic norm)’.} 본서의 중심 논제 중 하나는, 법체계의 토대들은
법적으로 무제한인 주권자에 대한 일반적 복종 습관(habit of obedience)에
있는 것이 아니라, 그 체계의 유효한 규칙들을 식별하는 권위 있는 기준들을
제공하는 \emphksans{궁극적 승인 규칙(ultimate rule of recognition)}에 있다는
것이다. 이 논제는 켈젠의 \emphksans{기초규범(basic norm)} 구상과 일정 부분
유사하며, 더 나아가서는 셀몬드(Salmond)의 충분히 전개되지 않은 ‘궁극적
법 원칙들(ultimate legal principles)’ 구상과도 가깝다(켈젠,
\emph{General Theory}, pp.~110--124, 131--134, 369--373, 395--396 참조;
셀몬드, \emph{Jurisprudence}, 11판, p.~137 및 부록 I 참조). 그러나
본서는 켈젠과는 다른 용어 체계를 채택하였으며, 그 이유는 아래의 주요한
점들에서 본서의 입장이 켈젠과 다르기 때문이다.

\begin{enumerate}
\def\labelenumi{\arabic{enumi}.}
\item
  승인 규칙이 존재하는지, 또 그것의 내용---즉 주어진 법체계에서의 유효성
  기준들(criteria of validity)---이 무엇인지는, 본서 전반에서
  경험적이면서도 복합적인 \emphksans{사실 문제(question of fact)}로 간주된다.
  이는 다음의 점이 사실이더라도 마찬가지이다: 즉, 법체계 안에서 활동하는
  법률가가 어떤 특정 규칙이 유효하다고 단언할(assert) 때, 그는 그 특정
  규칙의 유효성을 검사하기 위해 참조한 승인 규칙이 그 체계의 수용된 승인
  규칙으로서 존재한다는 사실을 \emphksans{명시적으로 진술하지는
  않지만(explicitly state)} \emphksans{암묵적으로 전제한다(tacitly
  presupposes)}. 만약 이 점에 대해 이의가 제기된다면, 이처럼
  전제되었지만 진술되지 않은 것은 사실들에 호소함으로써, 즉 적용해야 할
  법을 식별할 때 그 체계의 법원들과 공무담당자들이 보이는 실제 실천에
  호소함으로써 확립될 수 있다. 켈젠이 기초규범을 ‘법리학적 가설(juristic
  hypothesis)’(\emph{ibid.} xv), ‘가설적(hypothetical)’(\emph{ibid.}
  396), ‘상정된 궁극 규칙(postulated ultimate rule)’(\emph{ibid.} 113),
  ‘법리학적 의식 내에 존재하는 규칙(rule existing in the juristic
  consciousness)’(\emph{ibid.} 116), ‘가정(an assumption)’(\emph{ibid.}
  396) 등으로 분류하는 용어는, 이 책에서 강조되는 요점, 곧 “어떤
  법체계에서 법적 유효성의 기준들이 무엇인가”라는 물음이 사실 문제라는
  요점을 흐리게 하며, 실제로 그 요점과 양립하지 않을 수도 있다. 그것은
  사실 문제이지만, 하나의 규칙의 존재와 내용에 \emphksans{관한(about)}
  문제이다. Cf. 아고(Ago), “Positive Law and International Law,”
  \emph{American Journal of International Law} 제51권 (1957),
  pp.~703--707.
\item
  켈젠은 기초규범의 \emphksans{유효성(validity)}을 전제한다고 말한다. 그러나
  본문(pp.~108--110)에서 제시된 이유에 따르면, 일반적으로 수용된 승인
  규칙에 관해 그 존재 여부와는 별도로 그 유효성 또는
  비유효성(invalidity)을 묻는 것은 부적절하며, 그런 질문 자체가 성립할
  수 없다.
\item
  켈젠의 기초규범은 어떤 의미에서는 언제나 동일한 내용을 갖는다.
  왜냐하면 그것은 모든 법체계에서 단지 “헌법 또는 ‘최초의 헌법을 정한
  자들’에게 복종해야 한다”는 규칙이기 때문이다(\emph{General Theory},
  pp.~115--116). 그러나 이러한 일양성과 단순성의 외관은 오해를 낳을 수
  있다. 만약 하나의 헌법이 법의 다양한 원천들(sources of law)을
  명시하며, 그 법체계 내의 법원들과 공무담당자들이 실제로 그 헌법이
  제공하는 기준들에 따라 법을 식별한다는 의미에서 살아 있는 현실(living
  reality)이라면, 그러한 헌법은 수용된 것이며 실제로 존재하는 것이다.
  이런 경우, 헌법 또는 그것을 ‘정한’ 자들에게 복종해야 한다는 별도의
  규칙이 존재한다고 말하는 것은 불필요한 중복처럼 보인다. 특히, 영국처럼
  성문 헌법이 존재하지 않는 경우에는 더욱 그렇다. 이 경우, 법을 식별할
  때 특정한 유효성 기준들, 예컨대 의회 내 여왕에 의한 제정이 사용되어야
  한다는 규칙에 덧붙여 ‘헌법에 복종해야 한다’는 규칙이 자리할 곳은 없어
  보인다. 이 규칙 자체가 수용된 규칙이며, 이 규칙에 복종해야 한다는
  규칙을 말하는 것은 신비화할 뿐이다.
\item
  켈젠의 견해에 따르면(\emph{General Theory}, pp.~373--375, 408--410),
  특정한 법규칙을 유효한 것으로 간주하면서 동시에 그 법규칙이 요구하는
  행위를 금지하는 도덕규칙을 도덕적으로 구속력 있는 것으로 수용하는 것은
  논리적으로 불가능하다. 그러나 본서에서 제시한 법적 유효성에 대한
  설명에서는 그러한 결과가 전혀 도출되지 않는다. ‘\emphksans{기초(basic)}
  규범’이라는 표현 대신 ‘승인 규칙(rule of recognition)’이라는 표현을
  사용하는 이유 중 하나는, 법과 도덕 사이의 충돌에 관한 켈젠의 견해를
  따를 어떤 입장도 전제하지 않기 위함이다.
\end{enumerate}

\noteentry 101쪽. \notetitle{법의 원천들(sources of law).} 일부 저자들은 법들의
‘형식적(formal)’ 또는 ‘법적(legal)’ 원천들과 ‘역사적(historical)’ 또는
‘실질적(material)’ 원천들을 구분한다(Salmond, \emph{Jurisprudence},
11판, 제5장). Allen은 이 구분을 비판하지만(\emph{Law in the Making},
6판, p.~260), 이를 ‘원천(source)’이라는 단어의 두 가지 의미의 구분으로
해석하면 이 구분은 중요한 것이다(Kelsen, \emph{General Theory},
pp.~131--132, 152--153 참조). 한 의미에서(즉 ‘실질적(material)’,
‘역사적(historical)’ 의미에서) 원천이란, 특정한 시간과 장소에서 주어진
법규칙이 존재하게 된 원인적 또는 역사적 영향력들을 의미한다. 이 의미에서
현대 영국법의 특정 규칙들의 원천은 로마법 규칙일 수도 있고, 교회법
규칙이나 심지어 대중 도덕 규칙일 수도 있다. 그러나 ‘제정법(statute)’이
법의 원천이라고 말할 때, ‘원천(source)’이라는 단어는 단순한
역사적·인과적 영향력이 아니라, 문제의 법체계에서 수용된 법적 유효성의
기준들(criteria of legal validity) 중 하나를 가리킨다. 권능을 가진
입법부(competent legislature)에 의해 제정법으로 제정되었다는 점은,
주어진 제정법 규칙(statutory rule)이 단순히 존재하게 된
\emphksans{원인(cause)}이 아니라 그것이 유효한 법(valid law)인
\emphksans{이유(reason)}이다. 주어진 법규칙의 역사적 원인과 그 유효성의 이유
사이의 이러한 구분은, 입법부에 의한 제정, 관습적 실천, 선례 등을 유효한
법의 식별 표지들(identifying marks)로 수용하는 승인 규칙(rule of
recognition)이 체계 내에 존재할 때만 가능하다.

그러나 실제 실천에서는 역사적 또는 인과적 원천들과 법적 또는 형식적
원천들 간의 이 분명한 구분이 흐려질 수 있으며, Allen과 같은 저자들이 이
구분을 비판하게 된 것도 바로 이 점 때문이다. 어떤 체계에서
제정법(statute)이 법의 형식적 또는 법적 원천인 경우, 법원은 사건을
결정하면서 관련 제정법에 주의를 기울여야 할 구속을 받지만, 물론 제정법
문언의 의미를 해석하는 데 상당한 자유가 남겨져 있기는 하다(제7장 제1절
참조). 그러나 때때로 판사에게 남겨지는 것은 해석의 자유보다 훨씬 더
크다. 즉, 자기 앞의 사건을 결정하는 제정법이나 법의 다른 형식적 원천이
없다고 판단할 경우, 그는 예컨대 학설휘찬(Digest)의 한 문헌이나 프랑스
법학자의 저술에 자기 결정을 근거지을 수도 있다(예: Allen, \emph{op.
cit.}, p.~260 이하 참조). 해당 법체계가 그러한 원천들을 사용하라고
그에게 \emphksans{요구하는(require)} 것은 아니지만, 그가 그렇게 하는 것은
완전히 적절한 것으로 수용된다. 그러므로 이러한 저술들은 단순한 역사적
또는 인과적 영향력 이상이다. 그러한 저술들은 결정들을 위한 ‘좋은
이유들(good reasons)’로 인정되기 때문이다. 아마도 우리는 이러한
원천들을, 제정법과 같은 ‘구속적(mandatory)’ 법적 또는 형식적 원천들과
역사적 또는 실질적 원천들 양자와 구별하기 위해, ‘허용적(permissive) 법적
원천들’이라고 부를 수 있을 것이다.

\noteentry 103쪽. \notetitle{법적 유효성과 실효성(Legal validity and efficacy).} 켈젠은
전체적으로 실효적인 법질서의 실효성과 개별 규범의 실효성을
구분한다(\emph{General Theory}, pp.~41--2, 118--22). 그에게 있어 하나의
규범이 유효하다는 것은, 만약 그리고 오직 그 경우에만 (if, and only if),
그것이 \emphksans{전반적으로(on the whole)} 실효적인 체계에 속할 때이다. 그는
이 견해를 좀 더 난해하게 다음과 같이 표현하기도 한다. 즉, 체계 전체의
실효성은 그 체계의 규칙 유효성에 대한 \emphksans{필요조건(conditio sine qua
non)}이지, \emphksans{충분조건(conditio per quam)}은 아니라고
말한다(\emph{sed quaere}). 이 구분의 요점은 본서의 용어로 다음과 같이
정리될 수 있다. 법체계의 일반적인 실효성은 승인 규칙(rule of
recognition)이 제공하는 유효성 기준은 아니지만, 체계 내의 어떤 규칙이 그
유효성 기준들을 참조하여 그 체계의 유효한 규칙으로 식별될 때마다 언제나
암묵적으로 전제되며 명시적으로 진술되지는 않는다. 그리고 체계가
일반적으로 실효성을 갖추고 있지 않다면, 유효성에 관한 의미 있는 진술
자체가 불가능해진다. 본문의 입장은 이 점에서 켈젠과 다르다. 여기서는
체계의 실효성이 유효성 진술이 이루어지는 \emphksans{정상적 맥락(normal
context)}이긴 하지만, 특별한 상황에서는 그 체계가 더 이상 실효적이지
않더라도 유효성 진술이 여전히 의미를 가질 수 있다고 논변되기
때문이다(\emph{ante}, p.~104 참조).

켈젠은 또한 \emphksans{불사용에 의한 폐지(desuetudo)} 항목에서, 하나의
법체계가 특정 규칙의 유효성을 그 규칙의 지속적인 실효성에 의존하도록 할
가능성에 대해서도 논의한다. 이러한 경우, (개별 규칙의) 실효성은 체계의
유효성 기준 중 하나가 되며, 단순한 ‘전제(presupposition)’는
아니다(\emph{op. cit.}, pp.~119--22).

\noteentry 104쪽. \notetitle{유효성과 예측(Validity and prediction).} 법이 유효하다는
진술은 미래의 사법적 행위와 그에 따른 특유의 동기 부여 감정에 대한
예측이라는 견해에 대해서는 Ross, \emph{On Law and Justice}, 제1장 및
제2장 참조. 이에 대한 비판은 Hart, `Scandinavian Realism'
(\emph{Cambridge Law Journal}, 1959년)을 보라.

\noteentry 106쪽. \notetitle{개정권이 제한된 헌법(Constitutions with limited amending powers).} 서독과 터키의 사례는 제4장 주석(\emph{ante}, p.~290)을
참조하라.

\noteentry 111쪽. \notetitle{관행적 범주와 헌정 구조(Conventional categories and constitutional structures).} ‘법(law)’과 ‘관행(convention)’으로의 전부를
포괄한다고 주장되는 구분에 대해서는 Dicey, \emph{Law of the
Constitution}, 10판, pp.~23 이하; Wheare, \emph{Modern Constitutions},
제1장 참조.

\noteentry 111쪽. \notetitle{승인 규칙(rule of recognition): 법인가 사실인가?} 정치적
사실로 분류할 것인지에 대한 찬반 논증은 Wade, `The Basis of Legal
Sovereignty', \emph{Cambridge Law Journal} (1955), 특히 p.~189, 그리고
Marshall, \emph{Parliamentary Sovereignty and the Commonwealth},
pp.~43--46을 참조하라.

\noteentry 112쪽. \notetitle{법체계의 존재, 습관적 복종, 승인 규칙의 수용(The existence of a legal system, habitual obedience, and the acceptance of the rule of recognition).} 일반 시민의 복종과 공무담당자들 쪽에서의 헌정 규칙 수용을
모두 포함하는 복합적 사회 현상을 지나치게 단순화할 위험에 대해서는 제4장
제1절, pp.~60--61 및 Hughes, `The Existence of a Legal System', \emph{35
New York University LR} (1960), p.~1010 참조. Hughes는 이 점에서 Hart가
`Legal and Moral Obligation' (\emph{Essays in Moral Philosophy}, Melden
편, 1958)에서 사용한 용어를 정당하게 비판하고 있다.

\noteentry 118쪽. \notetitle{법질서의 부분적 붕괴(Partial breakdown of legal order).} 법체계가 완전히 정상적으로 존재하는 상태와 전혀 존재하지 않는 상태
사이의 많은 가능한 중간 상태들 중 극히 일부만이 본문에서 언급되었다.
혁명(revolution)에 대한 법적 관점에서의 논의는 Kelsen, \emph{General
Theory}, pp.~117 ff., 219 ff.를 참조하고, Cattaneo, \emph{Il Concetto di
Revoluzione nella Scienza del Diritto} (1960)에서는 상세히 다루어진다.
적군 점령(enemy occupation)에 의한 법체계의 중단은 다양한 형태를 띨 수
있으며, 이 중 일부는 국제법상으로 범주화되어 있다. 이에 대해서는 McNair,
`Municipal Effects of Belligerent Occupation', \emph{57 LQR} (1941),
그리고 Goodhart, `An Apology for Jurisprudence' (\emph{Interpretations
of Modern Legal Philosophies}, pp.~288 ff.)의 이론적 논의를 참조하라.

\noteentry 120쪽. \notetitle{법체계의 발생학(The embryology of a legal system).} Wheare,
\emph{The Statute of Westminster and Dominion Status,} 5판에서 추적되는
식민지에서 도미니언으로의 발전 과정은 법이론 연구에 보람 있는 분야이다.
또한 Latham, \emph{The Law and the Commonwealth} (1949)을 참조하라.
Latham은 영연방(Commonwealth)의 헌정 발전을 ‘지역적 뿌리(local root)’를
가진 새로운 기초규범의 성장이라는 관점에서 최초로 해석하였다. 또한
Marshall, 앞서 인용한 저서, 특히 캐나다에 대한 제7장을 참고하고, Wheare,
\emph{The Constitutional Structure of the Commonwealth} (1960), 제4장
‘자생성(Autochthony)’도 참고하라.

\noteentry 121쪽. \notetitle{입법 권한의 포기(Renunciation of legislative power).} 웨스트민스터 제정법(Statute of Westminster) 제4조의 법적 효력에 대한
논의는 Wheare, \emph{The Statute of Westminster and Dominion Status},
5th edn., pp.~297--8을 참조하라. 또한 \emph{British Coal Corporation} v.
\emph{The King} (1935), AC 500; Dixon, `The Law and the Constitution',
\emph{51 LQR} (1935); Marshall, 앞서 인용한 저서, pp.~146 ff.; 그리고
제7장 제4절도 함께 보라.

\noteentry 121쪽. \notetitle{모국 체계에 의해 인정되지 않은 독립(Independence not recognized by the parent system).} 아일랜드 자유국(Irish Free State)에
대한 논의는 Wheare, 앞서 인용한 저서 참조; \emph{Moore} v. \emph{AG for
the Irish Free State} (1935), AC 484; \emph{Ryan} v. \emph{Lennon}
(1935), IRR 170 참조.

\noteentry 121쪽. \notetitle{법체계의 존재에 관한 사실적 단언들과 법 진술들(Factual assertions and statements of law concerning the existence of a legal system).} Kelsen은 국내법과 국제법 사이의 가능한 관계들(‘국가법의 우위성
또는 국제법의 우위성’)에 대한 자신의 설명(op. cit., pp.~373--83)에서,
법체계가 존재한다는 진술은 반드시 한 법체계의 관점에서 다른 법체계에
대해, 그 다른 체계를 ‘유효한(valid)’ 것으로 그리고 자기 자신과 하나의
단일 체계를 형성하는 것으로 수용하면서 행해지는 법 진술(statement of
law)이어야 한다고 전제한다. 이에 반해 상식적 견해에서는 국내법과
국제법이 별개의 법체계들을 구성한다고 보며, 어떤 법체계(국가적이든
국제적이든)가 존재한다는 진술을 사실 진술(statement of fact)로 취급한다.
Kelsen에게는 이러한 견해가 용납될 수 없는
‘다원주의(pluralism)’이다(Kelsen, loc. cit.; Jones, `The “Pure” Theory
of International Law', \emph{16 BYBIL} 1935 참조). 또한 Hart, `Kelsen's
Doctrine of the Unity of Law', \emph{Ethics and Social Justice},
\emph{Contemporary Philosophical Thought}, vol.~4 (New York, 1970) 참조.

\noteentry 122쪽. \notetitle{남아프리카(South Africa).} 남아프리카 헌정 위기에서 배울 수
있는 중요한 법리학적 교훈을 보다 심층적으로 검토하려면, Marshall, 앞서
인용한 저서, 제11장을 참조하라.

\subsection{제6장 제3판
주석}\label{uxc81c6uxc7a5-uxc81c3uxd310-uxc8fcuxc11d}

\noteentry 100--1쪽. \notetitle{법의 원천들(Sources of law).} 법이
\emphksans{원천들(sources)}을 가진다는 관념의 이론적 중요성에 대해서는 Joseph
Raz, \emph{The Authority of Law}, chap.~3을 보라. 보통법(common law)과
관습(custom)을 법의 원천들로 다룬 논의는 Gerald J. Postema,
\emph{Bentham and the Common Law Tradition} (Oxford University Press,
1986) 및 John Gardner, `Some Types of Law', 그의 저서 \emph{Law as a
Leap of Faith} 제3장을 참조하라.

Hart의 승인 규칙(rule of recognition) 구상(conception)에 대한 비판은
Joseph Raz, \emph{The Concept of a Legal System}, chap.~8, 특히
pp.~197--200; Joseph Raz, \emph{The Authority of Law}, chap.~5, 특히
pp.~90--97; Ronald Dworkin, \emph{Taking Rights Seriously}, pp.~39--45,
62--68을 참조하라. 승인 규칙을 의미론적 교설로 해석한 그의 견해는 Ronald
Dworkin, \emph{Law's Empire}, pp.~31--35에서 볼 수 있다.

\noteentry 101--3쪽. \notetitle{승인 규칙과 법원들의 실천(The rule of recognition and the practice of the courts).} Hart에게서 궁극적 승인 규칙들(ultimate
rules of recognition)은 공무담당자들의 관습적 실천들 속에 내재한다.
그런데 공무담당자란 법에 의해 그렇게 식별되는 사람이라면, 순환성이
다가오는 것처럼 보일 수 있다. 즉, 어떤 것이 법인 것은 오직 그것이 승인
규칙에 의해 식별되는 경우뿐이고, 어떤 것이 승인 규칙인 것은 오직 그것이
공무담당자들에 의해 실천되는 경우뿐이며, 어떤 사람이 공무담당자인 것은
오직 그가 법에 의해 그렇게 권한을 부여받은 경우뿐이라는 식이다. 이
문제의 해소 방안들(resolutions)에 관해서는 Neil MacCormick, \emph{H. L.
A. Hart}, pp.~108--120; Michael Bayles, \emph{Hart's Legal Philosophy},
pp.~81--83을 참조하라.

\noteentry 103--4쪽. \notetitle{법적 유효성(Legal validity).} Hart는
‘유효성(validity)’을 주로 체계 소속(system membership)의 문제로 다룬다.
이에 대비되는 견해로는 Kelsen, \emph{Pure Theory of Law}, pp.~10--15,
193--195가 있다. 그는 “‘유효한(valid)’은 그것이 구속력을 가진다는 것,
곧 한 개인이 그 규범이 정한 방식으로 행위해야 한다는 것을
의미한다”(p.~193)고 말한다. 법적 맥락들에서 ‘유효성(validity)’의 다양한
의미들(senses)에 대해서는 J. W. Harris, \emph{Law and Legal Science: An
Inquiry into the Concepts Legal Rule and Legal System} (Oxford
University Press, 1979), chap.~4를 참조하라. Hart의 유효성 설명에 대한
평가는 Joseph Raz, \emph{The Authority of Law}, chap.~8에 실려 있다.
유효성의 정도들(degrees of validity)이라는 가능성에 대해서는 John
Finnis, \emph{Natural Law and Natural Rights}, pp.~276--281을 보라.

\noteentry 103--5쪽. \notetitle{법의 실효성(The efficacy of law).} Hart의 실효성 정의는
의무부과 규칙들(mandatory rules)에만 적용된다(“특정한 행위를 요구하는
법규칙이 대체로 복종된다는 사실”: p.~103). 그렇다면 허용적
규칙들(permissive rules) 혹은 권한-부여 규칙들(power-conferring rules)의
실효성은 어떻게 되는가? 이에 대해서는 Joseph Raz, \emph{The Concept of a
Legal System}, chap.~9; \emph{The Authority of Law}, chap.~5, 특히
pp.~85--90을 참조하라. 또한 Gerald Postema, `Conformity, Custom, and
Congruence: Rethinking the Efficacy of Law', Matthew H. Kramer \emph{et
al}. 편, \emph{The Legacy of H. L. A. Hart: Legal, Political, and Moral
Philosophy} (Oxford University Press, 2008)도 함께 보라.

\noteentry 105--10쪽. \notetitle{궁극적이고 최고인 승인 규칙(The rule of recognition as ultimate and supreme).} Hart는 ‘승인 규칙(rule of recognition)’이라는
용어를 한 법체계 안의 다른 모든 규칙들의 유효성을 평가할 기준들을
제공하지만, 그 자체는 다른 어떤 규칙에 의해서도 유효화되지(validated)
않는 하나의 관습적 사회 규칙을 지칭하는 데 사용한다. 다음 사항에
유의하라.

\begin{enumerate}
\def\labelenumi{(\arabic{enumi})}
\item
  승인 규칙은 판사들과 기타 공무담당자들의 관습적 규칙이지, 법으로
  제정된 것이 아니며 실정법(positive law)이 아니다. 특히, 그것은 형식적
  헌법(formal constitution)이나 그 일부가 아니다. (그러한 헌법은
  \emphksans{그 자체로} 승인 규칙에 의해 유효화되어야(validated) 한다.) 이
  점에 대해 혼란이 매우 많다. 이에 대한 논의로는 Kent Greenawalt, `The
  Rule of Recognition and the Constitution' (1987) 85 \emph{Michigan Law
  Review} 621; Matthew Adler and Kenneth Himma 편, \emph{The Rule of
  Recognition and the U.S. Constitution} (Oxford University Press, 2009)
  수록 논문들; John Gardner, `Can there be a Written Constitution',
  Leslie Green and Brian Leiter 편, \emph{Oxford Studies in Philosophy
  of Law,} vol.~I (Oxford University Press, 2011) 및 그의 저서 \emph{Law
  as a Leap of Faith} (Oxford University Press, 2012) 제4장을 참조하라.
\item
  Hart는 승인 규칙이 “그것이 그 집단의 규칙임을 보여 주는 확정적 인정
  표지(conclusive affirmative indication)를 제공하는 특징들”을
  명시한다고 말하고(94쪽), 그것이 “일차 규칙들을 확정적으로 식별하기
  위한 규칙”(95쪽)이라고 말하지만, 이를 승인 규칙의 기준들이 완전하고
  확정적이라는 의미로 이해해서는 안 된다. 그는 그것을 명시적으로
  부정한다(147--154쪽, 257--259쪽 참조).
\item
  Hart는 각 법체계에 정확히 하나의 승인 규칙이 존재한다고 아무런 논증
  없이 전제한다. 그러나 이는 의문의 여지가 있다. Joseph Raz, \emph{The
  Authority of Law}, 95--96쪽 참조.
\end{enumerate}

\noteentry 114--16쪽. \notetitle{공무담당자들과 이차 규칙들(Officials and secondary rules).} 흔히 두 가지 서로 다른 질문이 융합된다: (a) 누구의 행위가 승인
규칙을 \emphksans{구성하는가}? (b) 누구의 행위가 승인 규칙에 의해
\emphksans{규율되는가}? (a)에 답할 때 Hart는 때로
‘공무담당자들(officials)’(115, 116쪽)을, 때로는 ‘판사들(judges)’(108,
116쪽, 256쪽 참조)을 언급한다. Joseph Raz는 판사와 같은 법 적용
공무담당자들이 특히 중요하다고 논변한다: \emph{The Authority of Law},
chap.~6 참조. (b)에 관해서는 의문의 여지가 없다. 승인 규칙은 모든
공무담당자를 구속한다. 공무담당자들의 우위성, 그리고 법 적용
공무담당자들의 우위성은 정치적 도덕의 교설이 아니라는 점에 유의하라.
법원들이 실제로 그러한 영향력을 갖는 것이 유감스러울 수 있다는 점도
Hart가 논변하는 모든 것과 양립한다. 이에 대한 정치적 문제는 예컨대 Mark
Tushnet, \emph{Taking the Constitution Away from the Courts} (Princeton
University Press, 1999)를 보라.

\noteentry 116--17쪽. \notetitle{복종과 일반 시민(Obedience and the ordinary citizen).} 이 부분들은 Hart가 법을 사회 제도로 평가하는 맥락을 이해하는 데 핵심이
되는 대목들이다. 그는 법이 도덕적 위험들을 가져온다는 점뿐 아니라, 그
위험들이 법의 본성과 긴밀하게 연결되어 있다는 점을 논변한다. 이에 대한
논의로는 Jeremy Waldron, `All We Like Sheep' (1999) 12 \emph{Canadian
Journal of Law and Jurisprudence} 169 및 Leslie Green, `Positivism and
the Inseparability of Law and Morals' (2008) 83 \emph{New York
University Law Review} 1035, 특히 1052--1054쪽을 참조하라.

\noteentry 117--23쪽. \notetitle{법체계의 병리학(Pathology of a legal system).} John
Finnis, `Revolutions and Continuity of Law', 그의 저서 \emph{Philosophy
of Law: Collected Essays Volume IV} (Oxford University Press, 2011)의
제21장도 보라. 법체계들 사이의 상호작용 가능성들은 Neil MacCormick,
\emph{Questioning Sovereignty: Law, State and Nation in the European
Commonwealth} (Oxford University Press, 1999), chap.~7에서 탐구된다.
}
